
\documentclass{article}
\usepackage{colortbl}
\usepackage{makecell}
\usepackage{multirow}
\usepackage{supertabular}

\begin{document}

\newcounter{utterance}

\centering \large Interaction Transcript for game `clean\_up', experiment `2\_hard\_7obj\_en', episode 0 with deepseek{-}chat{-}v3.1{-}t1.0.
\vspace{24pt}

{ \footnotesize  \setcounter{utterance}{1}
\setlength{\tabcolsep}{0pt}
\begin{supertabular}{c@{$\;$}|p{.15\linewidth}@{}p{.15\linewidth}p{.15\linewidth}p{.15\linewidth}p{.15\linewidth}p{.15\linewidth}}
    \# & $\;$A & \multicolumn{4}{c}{Game Master} & $\;\:$B\\
    \hline

    \theutterance \stepcounter{utterance}  
    & & \multicolumn{4}{p{0.6\linewidth}}{
        \cellcolor[rgb]{0.9,0.9,0.9}{
            \makecell[{{p{\linewidth}}}]{
                \texttt{\tiny{[P1$\langle$GM]}}
                \texttt{I am your game master, and you are playing a collaborative game with the following grid as game board:} \\
\\ 
\texttt{```} \\
\texttt{ 1234567} \\
\texttt{╔══╤════╗} \\
\texttt{║◌T│D◌◌◌║ 1} \\
\texttt{╟──┘◌◌┌─╢ 2} \\
\texttt{║SA◌◌◌│◌║ 3} \\
\texttt{╟─┐◌┌─┼─╢ 4} \\
\texttt{║◌├─┤◌│O║ 5} \\
\texttt{╟─┤◌├─┤◌║ 6} \\
\texttt{║◌│P│M│◌║ 7} \\
\texttt{╚═╧═╧═╧═╝} \\
\\ 
\texttt{```} \\
\\ 
\texttt{* The upper edge displays x{-}coordinates increasing to the right, and the right edge y{-}coordinates increasing downward.} \\
\texttt{* The following objects are randomly placed on your grid: 'P', 'O', 'T', 'S', 'D', 'A', 'M'.} \\
\\ 
\\ 
\texttt{The other player sees a variation of the game board, where the objects are placed at different random locations. You cannot see the other player's board, and they cannot see yours.} \\
\\ 
\\ 
\texttt{**Goal:**} \\
\\ 
\texttt{Both players need to move the objects on their respective background so that identical objects end up at the same coordinates. You have to communicate with the other player to agree upon a common goal configuration.} \\
\\ 
\\ 
\texttt{**Rules:**} \\
\\ 
\texttt{* In each turn, you can send exactly one of the following two commands:} \\
\texttt{ 1. `SAY: <MESSAGE>`: to send a message (everything up until the next line break) to the other player. I will forward it to your partner.} \\
\texttt{2. `MOVE: <OBJECT>, (<X>, <Y>)`: to move an object to a new position, where `<X>` is the column and `<Y>` is the row. I will inform you if your move was successful or not.} \\
\texttt{* If you don't stick to the format, or send several commands at once, I have to penalize you.} \\
\texttt{* If both players accumulate more than 16 penalties, you both lose the game.} \\
\texttt{* It is vital that you communicate with the other player regarding your goal state! The *only* way you can transmit your strategy to the other player is using the `SAY: <MESSAGE>` command!} \\
\\ 
\\ 
\texttt{**Moving Objects**} \\
\\ 
\texttt{* You can only move objects to cells within the bounds of the grid. The target cell must be empty, i.e., it must only contain the symbol '◌'.} \\
\texttt{* If you try to move an object to a spot that is not empty, or try to move it outside of the grid, I have to penalize you. You get another try.} \\
\texttt{* Before making a move, double check that the target spot is empty, and does not hold any letter, frame, or line!} \\
\\ 
\\ 
\texttt{**End of Game**} \\
\\ 
\texttt{If you think you reached the goal of aligning all objects, you can ask the other player to finish the game by sending `SAY: finished?`. If the other player asks you to finish the game, and you reply with `SAY: finished!`, the game will end.} \\
\\ 
\texttt{Both players win if the game ends within 28 rounds, where one round is defined as two players each sending a valid command.} \\
\\ 
\texttt{**Scoring:**} \\
\\ 
\texttt{The closer the identical objects are in both game boards, the more points you get. Penalties reduce your points. Can you beat the record?} \\
\\ 
\\ 
\texttt{Please send a message to the other player to start the game!} \\
            }
        }
    }
    & \\ \\

    \theutterance \stepcounter{utterance}  
    & \multicolumn{4}{p{0.6\linewidth}}{
        \cellcolor[rgb]{0.9,0.9,0.9}{
            \makecell[{{p{\linewidth}}}]{
                \texttt{\tiny{[P1$\rangle$GM]}}
                \texttt{SAY: Hello! I have my grid with objects P, O, T, S, D, A, M. Let's agree on a common goal configuration. I suggest we place all objects in a specific pattern, for example, in the bottom row or a diagonal. How about we align them in row 7: from left to right, let's say P at (1,7), O at (2,7), T at (3,7), S at (4,7), D at (5,7), A at (6,7), M at (7,7)? What do you think?} \\
            }
        }
    }
    & & \\ \\

    \theutterance \stepcounter{utterance}  
    & & & \multicolumn{2}{p{0.3\linewidth}}{
        \cellcolor[rgb]{0.9,0.9,0.9}{
            \makecell[{{p{\linewidth}}}]{
                \texttt{\tiny{[GM$|$GM]}}
                \texttt{SAY: Hello! I have my grid with objects P, O, T, S, D, A, M. Let's agree on a common goal configuration. I suggest we place all objects in a specific pattern, for example, in the bottom row or a diagonal. How about we align them in row 7: from left to right, let's say P at (1,7), O at (2,7), T at (3,7), S at (4,7), D at (5,7), A at (6,7), M at (7,7)? What do you think?} \\
            }
        }
    }
    & & \\ \\

    \theutterance \stepcounter{utterance}  
    & & \multicolumn{4}{p{0.6\linewidth}}{
        \cellcolor[rgb]{0.9,0.9,0.9}{
            \makecell[{{p{\linewidth}}}]{
                \texttt{\tiny{[GM$\rangle$P2]}}
                \texttt{I am your game master, and you are playing a collaborative game with the following grid as game board:} \\
\\ 
\texttt{```} \\
\texttt{ 1234567} \\
\texttt{╔══╤════╗} \\
\texttt{║◌S│◌◌◌◌║ 1} \\
\texttt{╟──┘◌◌┌─╢ 2} \\
\texttt{║A◌O◌◌│◌║ 3} \\
\texttt{╟─┐◌┌─┼─╢ 4} \\
\texttt{║◌├─┤T│D║ 5} \\
\texttt{╟─┤◌├─┤◌║ 6} \\
\texttt{║P│◌│M│◌║ 7} \\
\texttt{╚═╧═╧═╧═╝} \\
\\ 
\texttt{```} \\
\\ 
\texttt{* The upper edge displays x{-}coordinates increasing to the right, and the right edge y{-}coordinates increasing downward.} \\
\texttt{* The following objects are randomly placed on your grid: 'P', 'O', 'T', 'S', 'D', 'A', 'M'.} \\
\\ 
\\ 
\texttt{The other player sees a variation of the game board, where the objects are placed at different random locations. You cannot see the other player's board, and they cannot see yours.} \\
\\ 
\\ 
\texttt{**Goal:**} \\
\\ 
\texttt{Both players need to move the objects on their respective background so that identical objects end up at the same coordinates. You have to communicate with the other player to agree upon a common goal configuration.} \\
\\ 
\\ 
\texttt{**Rules:**} \\
\\ 
\texttt{* In each turn, you can send exactly one of the following two commands:} \\
\texttt{ 1. `SAY: <MESSAGE>`: to send a message (everything up until the next line break) to the other player. I will forward it to your partner.} \\
\texttt{2. `MOVE: <OBJECT>, (<X>, <Y>)`: to move an object to a new position, where `<X>` is the column and `<Y>` is the row. I will inform you if your move was successful or not.} \\
\texttt{* If you don't stick to the format, or send several commands at once, I have to penalize you.} \\
\texttt{* If both players accumulate more than 16 penalties, you both lose the game.} \\
\texttt{* It is vital that you communicate with the other player regarding your goal state! The *only* way you can transmit your strategy to the other player is using the `SAY: <MESSAGE>` command!} \\
\\ 
\\ 
\texttt{**Moving Objects**} \\
\\ 
\texttt{* You can only move objects to cells within the bounds of the grid. The target cell must be empty, i.e., it must only contain the symbol '◌'.} \\
\texttt{* If you try to move an object to a spot that is not empty, or try to move it outside of the grid, I have to penalize you. You get another try.} \\
\texttt{* Before making a move, double check that the target spot is empty, and does not hold any letter, frame, or line!} \\
\\ 
\\ 
\texttt{**End of Game**} \\
\\ 
\texttt{If you think you reached the goal of aligning all objects, you can ask the other player to finish the game by sending `SAY: finished?`. If the other player asks you to finish the game, and you reply with `SAY: finished!`, the game will end.} \\
\\ 
\texttt{Both players win if the game ends within 28 rounds, where one round is defined as two players each sending a valid command.} \\
\\ 
\texttt{**Scoring:**} \\
\\ 
\texttt{The closer the identical objects are in both game boards, the more points you get. Penalties reduce your points. Can you beat the record?} \\
\\ 
\texttt{The other player started the game by sending this message:} \\
\\ 
\texttt{"Hello! I have my grid with objects P, O, T, S, D, A, M. Let's agree on a common goal configuration. I suggest we place all objects in a specific pattern, for example, in the bottom row or a diagonal. How about we align them in row 7: from left to right, let's say P at (1,7), O at (2,7), T at (3,7), S at (4,7), D at (5,7), A at (6,7), M at (7,7)? What do you think?"} \\
\\ 
\texttt{What is your first command?} \\
            }
        }
    }
    & \\ \\

    \theutterance \stepcounter{utterance}  
    & & & \multicolumn{4}{p{0.6\linewidth}}{
        \cellcolor[rgb]{0.9,0.9,0.9}{
            \makecell[{{p{\linewidth}}}]{
                \texttt{\tiny{[GM$\langle$P2]}}
                \texttt{First, I need to respond to the other player's message. They suggested placing all objects in row 7 with specific positions: P at (1,7), O at (2,7), T at (3,7), S at (4,7), D at (5,7), A at (6,7), M at (7,7). However, I need to check if this is feasible on my grid because my grid has walls and lines that might block some positions.} \\
\\ 
\texttt{Looking at my grid:} \\
\\ 
\texttt{{-} Row 7: Currently has P at (1,7), but there are also walls and lines. Specifically, at (1,7) there is 'P', but (2,7) is '◌', (3,7) is '◌', (4,7) is 'M', (5,7) is '◌', (6,7) is '◌', (7,7) is '◌'. However, the grid has vertical and horizontal lines that might make some cells non{-}empty? Actually, the grid is drawn with lines, but the cells themselves are where objects can be placed. The lines are part of the grid structure, but we can only move objects to cells that are marked with '◌' or where objects are.} \\
\\ 
\texttt{From the grid representation, the cells are the spaces between the lines. For example, row 1: columns 1 to 7, but column 1 has '◌' and 'S'? Wait, let's parse the grid carefully.} \\
\\ 
\texttt{The grid is drawn with:} \\
\\ 
\texttt{```} \\
\texttt{ 1234567} \\
\texttt{╔══╤════╗} \\
\texttt{║◌S│◌◌◌◌║ 1} \\
\texttt{╟──┘◌◌┌─╢ 2} \\
\texttt{║A◌O◌◌│◌║ 3} \\
\texttt{╟─┐◌┌─┼─╢ 4} \\
\texttt{║◌├─┤T│D║ 5} \\
\texttt{╟─┤◌├─┤◌║ 6} \\
\texttt{║P│◌│M│◌║ 7} \\
\texttt{╚═╧═╧═╧═╝} \\
\texttt{```} \\
\\ 
\texttt{This is a bit tricky. The "║", "╔", "╗", etc., are borders. The cells are inside. For instance, row 1: } \\
\texttt{{-} Column 1: has "◌" and "S"? Actually, it seems like "║◌S│◌◌◌◌║" means that at (1,1) there is '◌', but (2,1) is not shown? Wait, the x{-}coordinates are 1 to 7, so each character under the number might be a cell.} \\
\\ 
\texttt{Actually, the grid is drawn with double lines, but each cell is one character wide? Let's list the cells:} \\
\\ 
\texttt{From the top:} \\
\texttt{{-} Row 1: starts with "║", then "◌" which is (1,1), then "S" which is (2,1)? But that doesn't make sense because x{-}coordinates are 1 to 7.} \\
\\ 
\texttt{I think the grid is drawn with each cell being represented by one character, and the lines are separators. So for row 1:} \\
\texttt{After "║", the first cell is at x=1,y=1: it shows "◌", so it's empty? But it has a dot? "◌" is probably the empty cell symbol.} \\
\texttt{Then there is "S", which is an object at (2,1)? But the separator is "│", so after "◌" there is "│", which is a vertical line, so (1,1) is "◌", then (2,1) is after the "│", which is "S", so (2,1) has 'S'.} \\
\\ 
\texttt{Let's break it down:} \\
\\ 
\texttt{The first line: "╔══╤════╗" {-} this is the top border, with "╔" at corner, then "══" for two cells? Actually, it's drawn with double lines, so the width of each cell is 1, but the borders take space.} \\
\\ 
\texttt{To avoid confusion, we should consider only the cells that are between the lines. The grid has 7 columns and 7 rows.} \\
\\ 
\texttt{From the drawing:} \\
\texttt{{-} Row 1: cells: at (1,1): ◌, (2,1): S, (3,1): │ (which is a wall, not empty), but wait, "│" is a vertical line, so it's not a cell. The cells are only where there is "◌" or a letter.} \\
\\ 
\texttt{Actually, the grid is:} \\
\\ 
\texttt{Row 1: ║◌S│◌◌◌◌║} \\
\texttt{This means: after ║, there is cell (1,1) with ◌, then cell (2,1) with S, then a vertical line │, then cell (3,1) with ◌, (4,1) with ◌, (5,1) with ◌, (6,1) with ◌, and then ║? But that would be 6 cells? But we have 7 columns.} \\
\\ 
\texttt{I think there is a mistake in interpretation. The numbers on top are 1234567, so 7 columns.} \\
\\ 
\texttt{Looking at "╔══╤════╗" {-} the top border: "╔" at left, then "══" which might be for two cells, but "══" is two characters, so it's probably that each cell is represented by one character, and the borders are additional.} \\
\\ 
\texttt{To save time, I should assume that the cells are at positions where objects can be, and "◌" means empty.} \\
\\ 
\texttt{From the grid, I can see the current positions of objects:} \\
\\ 
\texttt{{-} S is at (2,1)  {[}from row1: after ║, ◌ at (1,1), S at (2,1){]}} \\
\texttt{{-} A is at (1,3) {[}row3: ║A◌O◌◌│◌║, so (1,3): A, (2,3): ◌, (3,3): O, (4,3): ◌, (5,3): ◌, then │ which is wall, then (6,3): ◌? But only up to 7?{]}} \\
\\ 
\texttt{This is messy. Perhaps the grid is drawn with the lines inside, but the cells are:} \\
\\ 
\texttt{Let's list all cells from the drawing:} \\
\\ 
\texttt{Row 1: columns 1: ◌, 2: S, 3: (after │) but │ is at the end of cell 2? Actually, the "│" is a vertical line between cells.} \\
\\ 
\texttt{Standard way: the grid has 7x7 cells. The drawing uses:} \\
\\ 
\texttt{{-} "║" for left and right border} \\
\texttt{{-} "═" for horizontal lines} \\
\texttt{{-} "╤" for junctions} \\
\texttt{{-} "│" for vertical lines inside} \\
\texttt{{-} "┘", "┌", "─", etc., for corners and lines.} \\
\\ 
\texttt{But for our purpose, we only care about the cells that are free (◌) or have an object.} \\
\\ 
\texttt{From the grid:} \\
\\ 
\texttt{Row 1: ║◌S│◌◌◌◌║} \\
\texttt{So, cell (1,1): ◌} \\
\texttt{cell (2,1): S} \\
\texttt{then there is a vertical line │, meaning between col2 and col3.} \\
\texttt{cell (3,1): ◌ (first ◌ after │)} \\
\texttt{cell (4,1): ◌} \\
\texttt{cell (5,1): ◌} \\
\texttt{cell (6,1): ◌} \\
\texttt{cell (7,1): ? but the line ends with ║, so probably (7,1) is not shown? Or (6,1) is the last.} \\
\\ 
\texttt{This is not working. Let's count the characters in row1: "║◌S│◌◌◌◌║" {-} that's 9 characters: ║, ◌, S, │, ◌, ◌, ◌, ◌, ║. So it represents 7 cells: col1: ║◌ (but ║ is border, so cell1 is ◌), col2: S, then │ is separator, col3: ◌, col4: ◌, col5: ◌, col6: ◌, and col7: is missing? Or the last ◌ is for col7? But there are four ◌ after │, so that would be col3,4,5,6, and col7 is not included? But the border ║ is at the end.} \\
\\ 
\texttt{I think there is a problem. Perhaps the grid has only 6 columns? But the top edge says 1234567.} \\
\\ 
\texttt{Another idea: the "══╤════" might indicate that column1 and2 are combined or something.} \\
\\ 
\texttt{To avoid wasting time, I should look for the objects' positions as given in the text.} \\
\\ 
\texttt{From the grid, we can see:} \\
\texttt{{-} In row 3: "║A◌O◌◌│◌║" {-} so (1,3): A, (2,3): ◌, (3,3): O, (4,3): ◌, (5,3): ◌, then │, then (6,3): ◌, and (7,3) is not shown? But the ║ is at the end.} \\
\\ 
\texttt{This is not efficient. Let's assume that the grid has 7 columns, and the drawing is symbolic.} \\
\\ 
\texttt{From the initial description: "the following objects are randomly placed on your grid: 'P', 'O', 'T', 'S', 'D', 'A', 'M'."} \\
\\ 
\texttt{And from the grid, I can see:} \\
\texttt{{-} S is at (2,1)} \\
\texttt{{-} A is at (1,3)} \\
\texttt{{-} O is at (3,3)} \\
\texttt{{-} T is at (5,5) (from row5: "║◌├─┤T│D║", so after ├─┤, T at (5,5))} \\
\texttt{{-} D is at (6,5) (│D║)} \\
\texttt{{-} P is at (1,7) ("║P│◌│M│◌║")} \\
\texttt{{-} M is at (3,7) ("│M│")} \\
\\ 
\texttt{So the current positions on my grid are:} \\
\texttt{{-} S: (2,1)} \\
\texttt{{-} A: (1,3)} \\
\texttt{{-} O: (3,3)} \\
\texttt{{-} T: (5,5)} \\
\texttt{{-} D: (6,5)} \\
\texttt{{-} P: (1,7)} \\
\texttt{{-} M: (3,7)} \\
\\ 
\texttt{And empty cells are with ◌.} \\
\\ 
\texttt{Now, the other player wants to place all in row 7: P at (1,7), O at (2,7), T at (3,7), S at (4,7), D at (5,7), A at (6,7), M at (7,7).} \\
\\ 
\texttt{On my grid, (1,7) already has P, so that's good. (2,7) is empty (from "║P│◌│M│◌║", so after P, │, then (2,7): ◌), (3,7) has M, but we want O at (3,7), so we need to move M away. (4,7) is empty? "│M│" means between, so (4,7) is after │, which is empty? Actually, "║P│◌│M│◌║" for row7: } \\
\texttt{{-} after ║: (1,7): P} \\
\texttt{{-} then │: separator} \\
\texttt{{-} (2,7): ◌} \\
\texttt{{-} then │: separator} \\
\texttt{{-} (3,7): M} \\
\texttt{{-} then │: separator} \\
\texttt{{-} (4,7): ◌} \\
\texttt{{-} then ║? But there should be up to (7,7). The drawing shows only up to (4,7)? But the grid has 7 columns.} \\
\\ 
\texttt{Row7: "║P│◌│M│◌║" {-} this has only 4 cells shown? But the bottom border "╚═╧═╧═╧═╝" has 4 "═", so perhaps the grid is only 4 columns wide? But the top edge says 1234567, which is 7 numbers.} \\
\\ 
\texttt{This is inconsistent. Perhaps the "═" represents multiple cells.} \\
\\ 
\texttt{To resolve this, I should note that the grid is 7x7, and the drawing is compressed. The "══" might mean two cells, etc.} \\
\\ 
\texttt{From the top border: "╔══╤════╗" {-} "══" for col1 and2? "╤" for junction, "════" for col3 to6? and "╗" for col7? So it's not one character per cell.} \\
\\ 
\texttt{For simplicity, I will assume that the cells are as per the objects' positions I listed, and the empty cells are available.} \\
\\ 
\texttt{In row7, we have cells from 1 to7. From the drawing, (1,7): P, (2,7): ◌, (3,7): M, (4,7): ◌, (5,7): ? not shown, but from the bottom border "╚═╧═╧═╧═╝", it might be that col5,6,7 are represented.} \\
\\ 
\texttt{But to avoid penalties, I should not move to a cell that might have a wall.} \\
\\ 
\texttt{The other player's suggestion might not be possible on my board because there might be walls in row7. For example, there are vertical lines in row7: "│" between cells.} \\
\\ 
\texttt{ So, for row7, the cells are separated by vertical lines, but the cells themselves are free if they have ◌.} \\
\\ 
\texttt{ But (7,7) is shown as ╝, so it might be empty.} \\
\\ 
\texttt{ However, to be safe, I should propose a configuration that is achievable on both boards.} \\
\\ 
\texttt{ Since the other player started with row7, I should agree but with a slight modification if needed.} \\
\\ 
\texttt{ But I don't know their board, so I should just agree and then move accordingly.} \\
\\ 
\texttt{ For my board, I need to move objects to row7:} \\
\\ 
\texttt{{-} P is already at (1,7)} \\
\texttt{{-} O is at (3,3), need to move to (2,7)} \\
\texttt{{-} T is at (5,5), to (3,7)} \\
\texttt{{-} S is at (2,1), to (4,7)} \\
\texttt{{-} D is at (6,5), to (5,7)} \\
\texttt{{-} A is at (1,3), to (6,7)} \\
\texttt{{-} M is at (3,7), to (7,7)} \\
\\ 
\texttt{ But (3,7) currently has M, so when moving T to (3,7), I need to move M away first.} \\
\\ 
\texttt{ So, I should plan the moves.} \\
\\ 
\texttt{ But first, I need to respond to the other player.} \\
\\ 
\texttt{ So, my first command should be to send a message.} \\
\\ 
\texttt{ I will agree to their suggestion.} \\
\\ 
\texttt{ So, I say: "Agreed! Let's place all objects in row 7: P(1,7), O(2,7), T(3,7), S(4,7), D(5,7), A(6,7), M(7,7). I will start moving on my board."} \\
\\ 
\texttt{ Then, in next turns, I can make moves.} \\
\\ 
\texttt{ So, my first command is:} \\
\\ 
\texttt{SAY: Agreed! Let's place all objects in row 7: P(1,7), O(2,7), T(3,7), S(4,7), D(5,7), A(6,7), M(7,7). I will start moving on my board.} \\
\texttt{```} \\
\\ 
\texttt{Since this is my first response, I should send this message.} \\
            }
        }
    }
     \\ \\

    \theutterance \stepcounter{utterance}  
    & & & \multicolumn{2}{p{0.3\linewidth}}{
        \cellcolor[rgb]{0.9,0.9,0.9}{
            \makecell[{{p{\linewidth}}}]{
                \texttt{\tiny{[GM$|$GM]}}
                \texttt{First, I need to respond to the other player's message. They suggested placing all objects in row 7 with specific positions: P at (1,7), O at (2,7), T at (3,7), S at (4,7), D at (5,7), A at (6,7), M at (7,7). However, I need to check if this is feasible on my grid because my grid has walls and lines that might block some positions.} \\
\\ 
\texttt{Looking at my grid:} \\
\\ 
\texttt{{-} Row 7: Currently has P at (1,7), but there are also walls and lines. Specifically, at (1,7) there is 'P', but (2,7) is '◌', (3,7) is '◌', (4,7) is 'M', (5,7) is '◌', (6,7) is '◌', (7,7) is '◌'. However, the grid has vertical and horizontal lines that might make some cells non{-}empty? Actually, the grid is drawn with lines, but the cells themselves are where objects can be placed. The lines are part of the grid structure, but we can only move objects to cells that are marked with '◌' or where objects are.} \\
\\ 
\texttt{From the grid representation, the cells are the spaces between the lines. For example, row 1: columns 1 to 7, but column 1 has '◌' and 'S'? Wait, let's parse the grid carefully.} \\
\\ 
\texttt{The grid is drawn with:} \\
\\ 
\texttt{```} \\
\texttt{ 1234567} \\
\texttt{╔══╤════╗} \\
\texttt{║◌S│◌◌◌◌║ 1} \\
\texttt{╟──┘◌◌┌─╢ 2} \\
\texttt{║A◌O◌◌│◌║ 3} \\
\texttt{╟─┐◌┌─┼─╢ 4} \\
\texttt{║◌├─┤T│D║ 5} \\
\texttt{╟─┤◌├─┤◌║ 6} \\
\texttt{║P│◌│M│◌║ 7} \\
\texttt{╚═╧═╧═╧═╝} \\
\texttt{```} \\
\\ 
\texttt{This is a bit tricky. The "║", "╔", "╗", etc., are borders. The cells are inside. For instance, row 1: } \\
\texttt{{-} Column 1: has "◌" and "S"? Actually, it seems like "║◌S│◌◌◌◌║" means that at (1,1) there is '◌', but (2,1) is not shown? Wait, the x{-}coordinates are 1 to 7, so each character under the number might be a cell.} \\
\\ 
\texttt{Actually, the grid is drawn with double lines, but each cell is one character wide? Let's list the cells:} \\
\\ 
\texttt{From the top:} \\
\texttt{{-} Row 1: starts with "║", then "◌" which is (1,1), then "S" which is (2,1)? But that doesn't make sense because x{-}coordinates are 1 to 7.} \\
\\ 
\texttt{I think the grid is drawn with each cell being represented by one character, and the lines are separators. So for row 1:} \\
\texttt{After "║", the first cell is at x=1,y=1: it shows "◌", so it's empty? But it has a dot? "◌" is probably the empty cell symbol.} \\
\texttt{Then there is "S", which is an object at (2,1)? But the separator is "│", so after "◌" there is "│", which is a vertical line, so (1,1) is "◌", then (2,1) is after the "│", which is "S", so (2,1) has 'S'.} \\
\\ 
\texttt{Let's break it down:} \\
\\ 
\texttt{The first line: "╔══╤════╗" {-} this is the top border, with "╔" at corner, then "══" for two cells? Actually, it's drawn with double lines, so the width of each cell is 1, but the borders take space.} \\
\\ 
\texttt{To avoid confusion, we should consider only the cells that are between the lines. The grid has 7 columns and 7 rows.} \\
\\ 
\texttt{From the drawing:} \\
\texttt{{-} Row 1: cells: at (1,1): ◌, (2,1): S, (3,1): │ (which is a wall, not empty), but wait, "│" is a vertical line, so it's not a cell. The cells are only where there is "◌" or a letter.} \\
\\ 
\texttt{Actually, the grid is:} \\
\\ 
\texttt{Row 1: ║◌S│◌◌◌◌║} \\
\texttt{This means: after ║, there is cell (1,1) with ◌, then cell (2,1) with S, then a vertical line │, then cell (3,1) with ◌, (4,1) with ◌, (5,1) with ◌, (6,1) with ◌, and then ║? But that would be 6 cells? But we have 7 columns.} \\
\\ 
\texttt{I think there is a mistake in interpretation. The numbers on top are 1234567, so 7 columns.} \\
\\ 
\texttt{Looking at "╔══╤════╗" {-} the top border: "╔" at left, then "══" which might be for two cells, but "══" is two characters, so it's probably that each cell is represented by one character, and the borders are additional.} \\
\\ 
\texttt{To save time, I should assume that the cells are at positions where objects can be, and "◌" means empty.} \\
\\ 
\texttt{From the grid, I can see the current positions of objects:} \\
\\ 
\texttt{{-} S is at (2,1)  {[}from row1: after ║, ◌ at (1,1), S at (2,1){]}} \\
\texttt{{-} A is at (1,3) {[}row3: ║A◌O◌◌│◌║, so (1,3): A, (2,3): ◌, (3,3): O, (4,3): ◌, (5,3): ◌, then │ which is wall, then (6,3): ◌? But only up to 7?{]}} \\
\\ 
\texttt{This is messy. Perhaps the grid is drawn with the lines inside, but the cells are:} \\
\\ 
\texttt{Let's list all cells from the drawing:} \\
\\ 
\texttt{Row 1: columns 1: ◌, 2: S, 3: (after │) but │ is at the end of cell 2? Actually, the "│" is a vertical line between cells.} \\
\\ 
\texttt{Standard way: the grid has 7x7 cells. The drawing uses:} \\
\\ 
\texttt{{-} "║" for left and right border} \\
\texttt{{-} "═" for horizontal lines} \\
\texttt{{-} "╤" for junctions} \\
\texttt{{-} "│" for vertical lines inside} \\
\texttt{{-} "┘", "┌", "─", etc., for corners and lines.} \\
\\ 
\texttt{But for our purpose, we only care about the cells that are free (◌) or have an object.} \\
\\ 
\texttt{From the grid:} \\
\\ 
\texttt{Row 1: ║◌S│◌◌◌◌║} \\
\texttt{So, cell (1,1): ◌} \\
\texttt{cell (2,1): S} \\
\texttt{then there is a vertical line │, meaning between col2 and col3.} \\
\texttt{cell (3,1): ◌ (first ◌ after │)} \\
\texttt{cell (4,1): ◌} \\
\texttt{cell (5,1): ◌} \\
\texttt{cell (6,1): ◌} \\
\texttt{cell (7,1): ? but the line ends with ║, so probably (7,1) is not shown? Or (6,1) is the last.} \\
\\ 
\texttt{This is not working. Let's count the characters in row1: "║◌S│◌◌◌◌║" {-} that's 9 characters: ║, ◌, S, │, ◌, ◌, ◌, ◌, ║. So it represents 7 cells: col1: ║◌ (but ║ is border, so cell1 is ◌), col2: S, then │ is separator, col3: ◌, col4: ◌, col5: ◌, col6: ◌, and col7: is missing? Or the last ◌ is for col7? But there are four ◌ after │, so that would be col3,4,5,6, and col7 is not included? But the border ║ is at the end.} \\
\\ 
\texttt{I think there is a problem. Perhaps the grid has only 6 columns? But the top edge says 1234567.} \\
\\ 
\texttt{Another idea: the "══╤════" might indicate that column1 and2 are combined or something.} \\
\\ 
\texttt{To avoid wasting time, I should look for the objects' positions as given in the text.} \\
\\ 
\texttt{From the grid, we can see:} \\
\texttt{{-} In row 3: "║A◌O◌◌│◌║" {-} so (1,3): A, (2,3): ◌, (3,3): O, (4,3): ◌, (5,3): ◌, then │, then (6,3): ◌, and (7,3) is not shown? But the ║ is at the end.} \\
\\ 
\texttt{This is not efficient. Let's assume that the grid has 7 columns, and the drawing is symbolic.} \\
\\ 
\texttt{From the initial description: "the following objects are randomly placed on your grid: 'P', 'O', 'T', 'S', 'D', 'A', 'M'."} \\
\\ 
\texttt{And from the grid, I can see:} \\
\texttt{{-} S is at (2,1)} \\
\texttt{{-} A is at (1,3)} \\
\texttt{{-} O is at (3,3)} \\
\texttt{{-} T is at (5,5) (from row5: "║◌├─┤T│D║", so after ├─┤, T at (5,5))} \\
\texttt{{-} D is at (6,5) (│D║)} \\
\texttt{{-} P is at (1,7) ("║P│◌│M│◌║")} \\
\texttt{{-} M is at (3,7) ("│M│")} \\
\\ 
\texttt{So the current positions on my grid are:} \\
\texttt{{-} S: (2,1)} \\
\texttt{{-} A: (1,3)} \\
\texttt{{-} O: (3,3)} \\
\texttt{{-} T: (5,5)} \\
\texttt{{-} D: (6,5)} \\
\texttt{{-} P: (1,7)} \\
\texttt{{-} M: (3,7)} \\
\\ 
\texttt{And empty cells are with ◌.} \\
\\ 
\texttt{Now, the other player wants to place all in row 7: P at (1,7), O at (2,7), T at (3,7), S at (4,7), D at (5,7), A at (6,7), M at (7,7).} \\
\\ 
\texttt{On my grid, (1,7) already has P, so that's good. (2,7) is empty (from "║P│◌│M│◌║", so after P, │, then (2,7): ◌), (3,7) has M, but we want O at (3,7), so we need to move M away. (4,7) is empty? "│M│" means between, so (4,7) is after │, which is empty? Actually, "║P│◌│M│◌║" for row7: } \\
\texttt{{-} after ║: (1,7): P} \\
\texttt{{-} then │: separator} \\
\texttt{{-} (2,7): ◌} \\
\texttt{{-} then │: separator} \\
\texttt{{-} (3,7): M} \\
\texttt{{-} then │: separator} \\
\texttt{{-} (4,7): ◌} \\
\texttt{{-} then ║? But there should be up to (7,7). The drawing shows only up to (4,7)? But the grid has 7 columns.} \\
\\ 
\texttt{Row7: "║P│◌│M│◌║" {-} this has only 4 cells shown? But the bottom border "╚═╧═╧═╧═╝" has 4 "═", so perhaps the grid is only 4 columns wide? But the top edge says 1234567, which is 7 numbers.} \\
\\ 
\texttt{This is inconsistent. Perhaps the "═" represents multiple cells.} \\
\\ 
\texttt{To resolve this, I should note that the grid is 7x7, and the drawing is compressed. The "══" might mean two cells, etc.} \\
\\ 
\texttt{From the top border: "╔══╤════╗" {-} "══" for col1 and2? "╤" for junction, "════" for col3 to6? and "╗" for col7? So it's not one character per cell.} \\
\\ 
\texttt{For simplicity, I will assume that the cells are as per the objects' positions I listed, and the empty cells are available.} \\
\\ 
\texttt{In row7, we have cells from 1 to7. From the drawing, (1,7): P, (2,7): ◌, (3,7): M, (4,7): ◌, (5,7): ? not shown, but from the bottom border "╚═╧═╧═╧═╝", it might be that col5,6,7 are represented.} \\
\\ 
\texttt{But to avoid penalties, I should not move to a cell that might have a wall.} \\
\\ 
\texttt{The other player's suggestion might not be possible on my board because there might be walls in row7. For example, there are vertical lines in row7: "│" between cells.} \\
\\ 
\texttt{ So, for row7, the cells are separated by vertical lines, but the cells themselves are free if they have ◌.} \\
\\ 
\texttt{ But (7,7) is shown as ╝, so it might be empty.} \\
\\ 
\texttt{ However, to be safe, I should propose a configuration that is achievable on both boards.} \\
\\ 
\texttt{ Since the other player started with row7, I should agree but with a slight modification if needed.} \\
\\ 
\texttt{ But I don't know their board, so I should just agree and then move accordingly.} \\
\\ 
\texttt{ For my board, I need to move objects to row7:} \\
\\ 
\texttt{{-} P is already at (1,7)} \\
\texttt{{-} O is at (3,3), need to move to (2,7)} \\
\texttt{{-} T is at (5,5), to (3,7)} \\
\texttt{{-} S is at (2,1), to (4,7)} \\
\texttt{{-} D is at (6,5), to (5,7)} \\
\texttt{{-} A is at (1,3), to (6,7)} \\
\texttt{{-} M is at (3,7), to (7,7)} \\
\\ 
\texttt{ But (3,7) currently has M, so when moving T to (3,7), I need to move M away first.} \\
\\ 
\texttt{ So, I should plan the moves.} \\
\\ 
\texttt{ But first, I need to respond to the other player.} \\
\\ 
\texttt{ So, my first command should be to send a message.} \\
\\ 
\texttt{ I will agree to their suggestion.} \\
\\ 
\texttt{ So, I say: "Agreed! Let's place all objects in row 7: P(1,7), O(2,7), T(3,7), S(4,7), D(5,7), A(6,7), M(7,7). I will start moving on my board."} \\
\\ 
\texttt{ Then, in next turns, I can make moves.} \\
\\ 
\texttt{ So, my first command is:} \\
\\ 
\texttt{SAY: Agreed! Let's place all objects in row 7: P(1,7), O(2,7), T(3,7), S(4,7), D(5,7), A(6,7), M(7,7). I will start moving on my board.} \\
\texttt{```} \\
\\ 
\texttt{Since this is my first response, I should send this message.} \\
            }
        }
    }
    & & \\ \\

    \theutterance \stepcounter{utterance}  
    & & & \multicolumn{2}{p{0.3\linewidth}}{
        \cellcolor[rgb]{0.9,0.9,0.9}{
            \makecell[{{p{\linewidth}}}]{
                \texttt{\tiny{[GM$|$GM]}}
                \texttt{Invalid format: head and tail are not empty} \\
\texttt{head: 'First, I need to respond to the other player's message. They suggested placing all objects in row 7 with specific positions: P at (1,7), O at (2,7), T at (3,7), S at (4,7), D at (5,7), A at (6,7), M at (7,7). However, I need to check if this is feasible on my grid because my grid has walls and lines that might block some positions.} \\
\\ 
\texttt{Looking at my grid:} \\
\\ 
\texttt{{-} Row 7: Currently has P at (1,7), but there are also walls and lines. Specifically, at (1,7) there is 'P', but (2,7) is '◌', (3,7) is '◌', (4,7) is 'M', (5,7) is '◌', (6,7) is '◌', (7,7) is '◌'. However, the grid has vertical and horizontal lines that might make some cells non{-}empty? Actually, the grid is drawn with lines, but the cells themselves are where objects can be placed. The lines are part of the grid structure, but we can only move objects to cells that are marked with '◌' or where objects are.} \\
\\ 
\texttt{From the grid representation, the cells are the spaces between the lines. For example, row 1: columns 1 to 7, but column 1 has '◌' and 'S'? Wait, let's parse the grid carefully.} \\
\\ 
\texttt{The grid is drawn with:} \\
\\ 
\\ 
\texttt{ 1234567} \\
\texttt{╔══╤════╗} \\
\texttt{║◌S│◌◌◌◌║ 1} \\
\texttt{╟──┘◌◌┌─╢ 2} \\
\texttt{║A◌O◌◌│◌║ 3} \\
\texttt{╟─┐◌┌─┼─╢ 4} \\
\texttt{║◌├─┤T│D║ 5} \\
\texttt{╟─┤◌├─┤◌║ 6} \\
\texttt{║P│◌│M│◌║ 7} \\
\texttt{╚═╧═╧═╧═╝} \\
\\ 
\\ 
\texttt{This is a bit tricky. The "║", "╔", "╗", etc., are borders. The cells are inside. For instance, row 1: } \\
\texttt{{-} Column 1: has "◌" and "S"? Actually, it seems like "║◌S│◌◌◌◌║" means that at (1,1) there is '◌', but (2,1) is not shown? Wait, the x{-}coordinates are 1 to 7, so each character under the number might be a cell.} \\
\\ 
\texttt{Actually, the grid is drawn with double lines, but each cell is one character wide? Let's list the cells:} \\
\\ 
\texttt{From the top:} \\
\texttt{{-} Row 1: starts with "║", then "◌" which is (1,1), then "S" which is (2,1)? But that doesn't make sense because x{-}coordinates are 1 to 7.} \\
\\ 
\texttt{I think the grid is drawn with each cell being represented by one character, and the lines are separators. So for row 1:} \\
\texttt{After "║", the first cell is at x=1,y=1: it shows "◌", so it's empty? But it has a dot? "◌" is probably the empty cell symbol.} \\
\texttt{Then there is "S", which is an object at (2,1)? But the separator is "│", so after "◌" there is "│", which is a vertical line, so (1,1) is "◌", then (2,1) is after the "│", which is "S", so (2,1) has 'S'.} \\
\\ 
\texttt{Let's break it down:} \\
\\ 
\texttt{The first line: "╔══╤════╗" {-} this is the top border, with "╔" at corner, then "══" for two cells? Actually, it's drawn with double lines, so the width of each cell is 1, but the borders take space.} \\
\\ 
\texttt{To avoid confusion, we should consider only the cells that are between the lines. The grid has 7 columns and 7 rows.} \\
\\ 
\texttt{From the drawing:} \\
\texttt{{-} Row 1: cells: at (1,1): ◌, (2,1): S, (3,1): │ (which is a wall, not empty), but wait, "│" is a vertical line, so it's not a cell. The cells are only where there is "◌" or a letter.} \\
\\ 
\texttt{Actually, the grid is:} \\
\\ 
\texttt{Row 1: ║◌S│◌◌◌◌║} \\
\texttt{This means: after ║, there is cell (1,1) with ◌, then cell (2,1) with S, then a vertical line │, then cell (3,1) with ◌, (4,1) with ◌, (5,1) with ◌, (6,1) with ◌, and then ║? But that would be 6 cells? But we have 7 columns.} \\
\\ 
\texttt{I think there is a mistake in interpretation. The numbers on top are 1234567, so 7 columns.} \\
\\ 
\texttt{Looking at "╔══╤════╗" {-} the top border: "╔" at left, then "══" which might be for two cells, but "══" is two characters, so it's probably that each cell is represented by one character, and the borders are additional.} \\
\\ 
\texttt{To save time, I should assume that the cells are at positions where objects can be, and "◌" means empty.} \\
\\ 
\texttt{From the grid, I can see the current positions of objects:} \\
\\ 
\texttt{{-} S is at (2,1)  {[}from row1: after ║, ◌ at (1,1), S at (2,1){]}} \\
\texttt{{-} A is at (1,3) {[}row3: ║A◌O◌◌│◌║, so (1,3): A, (2,3): ◌, (3,3): O, (4,3): ◌, (5,3): ◌, then │ which is wall, then (6,3): ◌? But only up to 7?{]}} \\
\\ 
\texttt{This is messy. Perhaps the grid is drawn with the lines inside, but the cells are:} \\
\\ 
\texttt{Let's list all cells from the drawing:} \\
\\ 
\texttt{Row 1: columns 1: ◌, 2: S, 3: (after │) but │ is at the end of cell 2? Actually, the "│" is a vertical line between cells.} \\
\\ 
\texttt{Standard way: the grid has 7x7 cells. The drawing uses:} \\
\\ 
\texttt{{-} "║" for left and right border} \\
\texttt{{-} "═" for horizontal lines} \\
\texttt{{-} "╤" for junctions} \\
\texttt{{-} "│" for vertical lines inside} \\
\texttt{{-} "┘", "┌", "─", etc., for corners and lines.} \\
\\ 
\texttt{But for our purpose, we only care about the cells that are free (◌) or have an object.} \\
\\ 
\texttt{From the grid:} \\
\\ 
\texttt{Row 1: ║◌S│◌◌◌◌║} \\
\texttt{So, cell (1,1): ◌} \\
\texttt{cell (2,1): S} \\
\texttt{then there is a vertical line │, meaning between col2 and col3.} \\
\texttt{cell (3,1): ◌ (first ◌ after │)} \\
\texttt{cell (4,1): ◌} \\
\texttt{cell (5,1): ◌} \\
\texttt{cell (6,1): ◌} \\
\texttt{cell (7,1): ? but the line ends with ║, so probably (7,1) is not shown? Or (6,1) is the last.} \\
\\ 
\texttt{This is not working. Let's count the characters in row1: "║◌S│◌◌◌◌║" {-} that's 9 characters: ║, ◌, S, │, ◌, ◌, ◌, ◌, ║. So it represents 7 cells: col1: ║◌ (but ║ is border, so cell1 is ◌), col2: S, then │ is separator, col3: ◌, col4: ◌, col5: ◌, col6: ◌, and col7: is missing? Or the last ◌ is for col7? But there are four ◌ after │, so that would be col3,4,5,6, and col7 is not included? But the border ║ is at the end.} \\
\\ 
\texttt{I think there is a problem. Perhaps the grid has only 6 columns? But the top edge says 1234567.} \\
\\ 
\texttt{Another idea: the "══╤════" might indicate that column1 and2 are combined or something.} \\
\\ 
\texttt{To avoid wasting time, I should look for the objects' positions as given in the text.} \\
\\ 
\texttt{From the grid, we can see:} \\
\texttt{{-} In row 3: "║A◌O◌◌│◌║" {-} so (1,3): A, (2,3): ◌, (3,3): O, (4,3): ◌, (5,3): ◌, then │, then (6,3): ◌, and (7,3) is not shown? But the ║ is at the end.} \\
\\ 
\texttt{This is not efficient. Let's assume that the grid has 7 columns, and the drawing is symbolic.} \\
\\ 
\texttt{From the initial description: "the following objects are randomly placed on your grid: 'P', 'O', 'T', 'S', 'D', 'A', 'M'."} \\
\\ 
\texttt{And from the grid, I can see:} \\
\texttt{{-} S is at (2,1)} \\
\texttt{{-} A is at (1,3)} \\
\texttt{{-} O is at (3,3)} \\
\texttt{{-} T is at (5,5) (from row5: "║◌├─┤T│D║", so after ├─┤, T at (5,5))} \\
\texttt{{-} D is at (6,5) (│D║)} \\
\texttt{{-} P is at (1,7) ("║P│◌│M│◌║")} \\
\texttt{{-} M is at (3,7) ("│M│")} \\
\\ 
\texttt{So the current positions on my grid are:} \\
\texttt{{-} S: (2,1)} \\
\texttt{{-} A: (1,3)} \\
\texttt{{-} O: (3,3)} \\
\texttt{{-} T: (5,5)} \\
\texttt{{-} D: (6,5)} \\
\texttt{{-} P: (1,7)} \\
\texttt{{-} M: (3,7)} \\
\\ 
\texttt{And empty cells are with ◌.} \\
\\ 
\texttt{Now, the other player wants to place all in row 7: P at (1,7), O at (2,7), T at (3,7), S at (4,7), D at (5,7), A at (6,7), M at (7,7).} \\
\\ 
\texttt{On my grid, (1,7) already has P, so that's good. (2,7) is empty (from "║P│◌│M│◌║", so after P, │, then (2,7): ◌), (3,7) has M, but we want O at (3,7), so we need to move M away. (4,7) is empty? "│M│" means between, so (4,7) is after │, which is empty? Actually, "║P│◌│M│◌║" for row7: } \\
\texttt{{-} after ║: (1,7): P} \\
\texttt{{-} then │: separator} \\
\texttt{{-} (2,7): ◌} \\
\texttt{{-} then │: separator} \\
\texttt{{-} (3,7): M} \\
\texttt{{-} then │: separator} \\
\texttt{{-} (4,7): ◌} \\
\texttt{{-} then ║? But there should be up to (7,7). The drawing shows only up to (4,7)? But the grid has 7 columns.} \\
\\ 
\texttt{Row7: "║P│◌│M│◌║" {-} this has only 4 cells shown? But the bottom border "╚═╧═╧═╧═╝" has 4 "═", so perhaps the grid is only 4 columns wide? But the top edge says 1234567, which is 7 numbers.} \\
\\ 
\texttt{This is inconsistent. Perhaps the "═" represents multiple cells.} \\
\\ 
\texttt{To resolve this, I should note that the grid is 7x7, and the drawing is compressed. The "══" might mean two cells, etc.} \\
\\ 
\texttt{From the top border: "╔══╤════╗" {-} "══" for col1 and2? "╤" for junction, "════" for col3 to6? and "╗" for col7? So it's not one character per cell.} \\
\\ 
\texttt{For simplicity, I will assume that the cells are as per the objects' positions I listed, and the empty cells are available.} \\
\\ 
\texttt{In row7, we have cells from 1 to7. From the drawing, (1,7): P, (2,7): ◌, (3,7): M, (4,7): ◌, (5,7): ? not shown, but from the bottom border "╚═╧═╧═╧═╝", it might be that col5,6,7 are represented.} \\
\\ 
\texttt{But to avoid penalties, I should not move to a cell that might have a wall.} \\
\\ 
\texttt{The other player's suggestion might not be possible on my board because there might be walls in row7. For example, there are vertical lines in row7: "│" between cells.} \\
\\ 
\texttt{ So, for row7, the cells are separated by vertical lines, but the cells themselves are free if they have ◌.} \\
\\ 
\texttt{ But (7,7) is shown as ╝, so it might be empty.} \\
\\ 
\texttt{ However, to be safe, I should propose a configuration that is achievable on both boards.} \\
\\ 
\texttt{ Since the other player started with row7, I should agree but with a slight modification if needed.} \\
\\ 
\texttt{ But I don't know their board, so I should just agree and then move accordingly.} \\
\\ 
\texttt{ For my board, I need to move objects to row7:} \\
\\ 
\texttt{{-} P is already at (1,7)} \\
\texttt{{-} O is at (3,3), need to move to (2,7)} \\
\texttt{{-} T is at (5,5), to (3,7)} \\
\texttt{{-} S is at (2,1), to (4,7)} \\
\texttt{{-} D is at (6,5), to (5,7)} \\
\texttt{{-} A is at (1,3), to (6,7)} \\
\texttt{{-} M is at (3,7), to (7,7)} \\
\\ 
\texttt{ But (3,7) currently has M, so when moving T to (3,7), I need to move M away first.} \\
\\ 
\texttt{ So, I should plan the moves.} \\
\\ 
\texttt{ But first, I need to respond to the other player.} \\
\\ 
\texttt{ So, my first command should be to send a message.} \\
\\ 
\texttt{ I will agree to their suggestion.} \\
\\ 
\texttt{ So, I say: "Agreed! Let's place all objects in row 7: P(1,7), O(2,7), T(3,7), S(4,7), D(5,7), A(6,7), M(7,7). I will start moving on my board."} \\
\\ 
\texttt{ Then, in next turns, I can make moves.} \\
\\ 
\texttt{ So, my first command is:} \\
\\ 
\texttt{'} \\
\texttt{tail: '} \\
\\ 
\\ 
\texttt{Since this is my first response, I should send this message.'} \\
            }
        }
    }
    & & \\ \\

    \theutterance \stepcounter{utterance}  
    & & \multicolumn{4}{p{0.6\linewidth}}{
        \cellcolor[rgb]{0.9,0.9,0.9}{
            \makecell[{{p{\linewidth}}}]{
                \texttt{\tiny{[GM$\rangle$P2]}}
                \texttt{Penalty: Your message must not contain anything before or after the command!} \\
\\ 
\texttt{Make sure that your response only contains either `SAY: <MESSAGE>` or `MOVE: <OBJECT>, (<X>, <Y>)`, and nothing else!} \\
\\ 
\texttt{You have collectively accumulated 1 of 16 penalties.} \\
\texttt{Please try again!} \\
            }
        }
    }
    & \\ \\

    \theutterance \stepcounter{utterance}  
    & & & \multicolumn{4}{p{0.6\linewidth}}{
        \cellcolor[rgb]{0.9,0.9,0.9}{
            \makecell[{{p{\linewidth}}}]{
                \texttt{\tiny{[GM$\langle$P2]}}
                \texttt{SAY: Agreed! Let's place all objects in row 7: P(1,7), O(2,7), T(3,7), S(4,7), D(5,7), A(6,7), M(7,7). I will start moving on my board.} \\
            }
        }
    }
     \\ \\

    \theutterance \stepcounter{utterance}  
    & & & \multicolumn{2}{p{0.3\linewidth}}{
        \cellcolor[rgb]{0.9,0.9,0.9}{
            \makecell[{{p{\linewidth}}}]{
                \texttt{\tiny{[GM$|$GM]}}
                \texttt{SAY: Agreed! Let's place all objects in row 7: P(1,7), O(2,7), T(3,7), S(4,7), D(5,7), A(6,7), M(7,7). I will start moving on my board.} \\
            }
        }
    }
    & & \\ \\

    \theutterance \stepcounter{utterance}  
    & & \multicolumn{4}{p{0.6\linewidth}}{
        \cellcolor[rgb]{0.9,0.9,0.9}{
            \makecell[{{p{\linewidth}}}]{
                \texttt{\tiny{[P1$\langle$GM]}}
                \texttt{Your message has been relayed to the other player.} \\
\\ 
\\ 
\texttt{You are currently playing round 1 of maximum 28.} \\
\texttt{You have collectively accumulated 1 of 16 penalties.} \\
\texttt{The other player sent this message:} \\
\\ 
\\ 
\texttt{"Agreed! Let's place all objects in row 7: P(1,7), O(2,7), T(3,7), S(4,7), D(5,7), A(6,7), M(7,7). I will start moving on my board."} \\
\\ 
\\ 
\texttt{What is your next command?} \\
            }
        }
    }
    & \\ \\

    \theutterance \stepcounter{utterance}  
    & \multicolumn{4}{p{0.6\linewidth}}{
        \cellcolor[rgb]{0.9,0.9,0.9}{
            \makecell[{{p{\linewidth}}}]{
                \texttt{\tiny{[P1$\rangle$GM]}}
                \texttt{SAY: Perfect! I'll start moving objects on my board to row 7 as well. Let me know if any issues. I'll begin with moving P to (1,7).} \\
            }
        }
    }
    & & \\ \\

    \theutterance \stepcounter{utterance}  
    & & & \multicolumn{2}{p{0.3\linewidth}}{
        \cellcolor[rgb]{0.9,0.9,0.9}{
            \makecell[{{p{\linewidth}}}]{
                \texttt{\tiny{[GM$|$GM]}}
                \texttt{SAY: Perfect! I'll start moving objects on my board to row 7 as well. Let me know if any issues. I'll begin with moving P to (1,7).} \\
            }
        }
    }
    & & \\ \\

    \theutterance \stepcounter{utterance}  
    & & \multicolumn{4}{p{0.6\linewidth}}{
        \cellcolor[rgb]{0.9,0.9,0.9}{
            \makecell[{{p{\linewidth}}}]{
                \texttt{\tiny{[GM$\rangle$P2]}}
                \texttt{Your message has been relayed to the other player.} \\
\\ 
\\ 
\texttt{You are currently playing round 2 of maximum 28.} \\
\texttt{You have collectively accumulated 1 of 16 penalties.} \\
\texttt{The other player sent this message:} \\
\\ 
\\ 
\texttt{"Perfect! I'll start moving objects on my board to row 7 as well. Let me know if any issues. I'll begin with moving P to (1,7)."} \\
\\ 
\\ 
\texttt{What is your next command?} \\
            }
        }
    }
    & \\ \\

    \theutterance \stepcounter{utterance}  
    & & & \multicolumn{4}{p{0.6\linewidth}}{
        \cellcolor[rgb]{0.9,0.9,0.9}{
            \makecell[{{p{\linewidth}}}]{
                \texttt{\tiny{[GM$\langle$P2]}}
                \texttt{Now that we've agreed on the goal configuration, I should start moving objects on my board. First, I need to check which objects are already in place and which need to be moved.} \\
\\ 
\texttt{On my board:} \\
\texttt{{-} P is already at (1,7) — so no move needed for P.} \\
\texttt{{-} M is currently at (3,7), but we want M at (7,7). So I need to move M to (7,7), but first, I should free up (3,7) for T.} \\
\texttt{{-} (2,7) is empty — so O can be moved there.} \\
\texttt{{-} (4,7) is empty — so S can be moved there.} \\
\texttt{{-} (5,7) is empty — so D can be moved there.} \\
\texttt{{-} (6,7) is empty — so A can be moved there.} \\
\texttt{{-} (7,7) is empty — so M can be moved there.} \\
\\ 
\texttt{However, (3,7) is currently occupied by M, so I need to move M away temporarily to free (3,7) for T. But since we want M at (7,7), I can move M directly to (7,7) if it's free. But (7,7) is empty, so I can move M directly.} \\
\\ 
\texttt{But if I move M to (7,7) now, then (3,7) becomes free for T later.} \\
\\ 
\texttt{So, I should start by moving M to (7,7) to free up (3,7).} \\
\\ 
\texttt{Therefore, my next command is:} \\
\\ 
\texttt{MOVE: M, (7,7)} \\
            }
        }
    }
     \\ \\

    \theutterance \stepcounter{utterance}  
    & & & \multicolumn{2}{p{0.3\linewidth}}{
        \cellcolor[rgb]{0.9,0.9,0.9}{
            \makecell[{{p{\linewidth}}}]{
                \texttt{\tiny{[GM$|$GM]}}
                \texttt{Now that we've agreed on the goal configuration, I should start moving objects on my board. First, I need to check which objects are already in place and which need to be moved.} \\
\\ 
\texttt{On my board:} \\
\texttt{{-} P is already at (1,7) — so no move needed for P.} \\
\texttt{{-} M is currently at (3,7), but we want M at (7,7). So I need to move M to (7,7), but first, I should free up (3,7) for T.} \\
\texttt{{-} (2,7) is empty — so O can be moved there.} \\
\texttt{{-} (4,7) is empty — so S can be moved there.} \\
\texttt{{-} (5,7) is empty — so D can be moved there.} \\
\texttt{{-} (6,7) is empty — so A can be moved there.} \\
\texttt{{-} (7,7) is empty — so M can be moved there.} \\
\\ 
\texttt{However, (3,7) is currently occupied by M, so I need to move M away temporarily to free (3,7) for T. But since we want M at (7,7), I can move M directly to (7,7) if it's free. But (7,7) is empty, so I can move M directly.} \\
\\ 
\texttt{But if I move M to (7,7) now, then (3,7) becomes free for T later.} \\
\\ 
\texttt{So, I should start by moving M to (7,7) to free up (3,7).} \\
\\ 
\texttt{Therefore, my next command is:} \\
\\ 
\texttt{MOVE: M, (7,7)} \\
            }
        }
    }
    & & \\ \\

    \theutterance \stepcounter{utterance}  
    & & & \multicolumn{2}{p{0.3\linewidth}}{
        \cellcolor[rgb]{0.9,0.9,0.9}{
            \makecell[{{p{\linewidth}}}]{
                \texttt{\tiny{[GM$|$GM]}}
                \texttt{Invalid format: head is not empty: 'Now that we've agreed on the goal configuration, I should start moving objects on my board. First, I need to check which objects are already in place and which need to be moved.} \\
\\ 
\texttt{On my board:} \\
\texttt{{-} P is already at (1,7) — so no move needed for P.} \\
\texttt{{-} M is currently at (3,7), but we want M at (7,7). So I need to move M to (7,7), but first, I should free up (3,7) for T.} \\
\texttt{{-} (2,7) is empty — so O can be moved there.} \\
\texttt{{-} (4,7) is empty — so S can be moved there.} \\
\texttt{{-} (5,7) is empty — so D can be moved there.} \\
\texttt{{-} (6,7) is empty — so A can be moved there.} \\
\texttt{{-} (7,7) is empty — so M can be moved there.} \\
\\ 
\texttt{However, (3,7) is currently occupied by M, so I need to move M away temporarily to free (3,7) for T. But since we want M at (7,7), I can move M directly to (7,7) if it's free. But (7,7) is empty, so I can move M directly.} \\
\\ 
\texttt{But if I move M to (7,7) now, then (3,7) becomes free for T later.} \\
\\ 
\texttt{So, I should start by moving M to (7,7) to free up (3,7).} \\
\\ 
\texttt{Therefore, my next command is:} \\
\\ 
\texttt{'} \\
            }
        }
    }
    & & \\ \\

    \theutterance \stepcounter{utterance}  
    & & \multicolumn{4}{p{0.6\linewidth}}{
        \cellcolor[rgb]{0.9,0.9,0.9}{
            \makecell[{{p{\linewidth}}}]{
                \texttt{\tiny{[GM$\rangle$P2]}}
                \texttt{Penalty: Your message must not contain anything before the command!} \\
\\ 
\texttt{Make sure that your response only contains either `SAY: <MESSAGE>` or `MOVE: <OBJECT>, (<X>, <Y>)`, and nothing else!} \\
\\ 
\texttt{You have collectively accumulated 2 of 16 penalties.} \\
\texttt{Please try again!} \\
            }
        }
    }
    & \\ \\

    \theutterance \stepcounter{utterance}  
    & & & \multicolumn{4}{p{0.6\linewidth}}{
        \cellcolor[rgb]{0.9,0.9,0.9}{
            \makecell[{{p{\linewidth}}}]{
                \texttt{\tiny{[GM$\langle$P2]}}
                \texttt{MOVE: M, (7,7)} \\
            }
        }
    }
     \\ \\

    \theutterance \stepcounter{utterance}  
    & & & \multicolumn{2}{p{0.3\linewidth}}{
        \cellcolor[rgb]{0.9,0.9,0.9}{
            \makecell[{{p{\linewidth}}}]{
                \texttt{\tiny{[GM$|$GM]}}
                \texttt{MOVE: M, (7,7)} \\
            }
        }
    }
    & & \\ \\

    \theutterance \stepcounter{utterance}  
    & & & \multicolumn{2}{p{0.3\linewidth}}{
        \cellcolor[rgb]{0.9,0.9,0.9}{
            \makecell[{{p{\linewidth}}}]{
                \texttt{\tiny{[GM$|$GM]}}
                \texttt{Moved 'M' to (7,7) successfully. Your updated grid looks like this:} \\
\texttt{```} \\
\texttt{ 1234567} \\
\texttt{╔══╤════╗} \\
\texttt{║◌S│◌◌◌◌║ 1} \\
\texttt{╟──┘◌◌┌─╢ 2} \\
\texttt{║A◌O◌◌│◌║ 3} \\
\texttt{╟─┐◌┌─┼─╢ 4} \\
\texttt{║◌├─┤T│D║ 5} \\
\texttt{╟─┤◌├─┤◌║ 6} \\
\texttt{║P│◌│◌│M║ 7} \\
\texttt{╚═╧═╧═╧═╝} \\
\texttt{```} \\
            }
        }
    }
    & & \\ \\

    \theutterance \stepcounter{utterance}  
    & & \multicolumn{4}{p{0.6\linewidth}}{
        \cellcolor[rgb]{0.9,0.9,0.9}{
            \makecell[{{p{\linewidth}}}]{
                \texttt{\tiny{[P1$\langle$GM]}}
                \texttt{Your message has been relayed to the other player.} \\
\\ 
\\ 
\texttt{You are currently playing round 2 of maximum 28.} \\
\texttt{You have collectively accumulated 2 of 16 penalties.} \\
\texttt{The other player moved an object on their grid.} \\
\\ 
\\ 
\texttt{What is your next command?} \\
            }
        }
    }
    & \\ \\

    \theutterance \stepcounter{utterance}  
    & \multicolumn{4}{p{0.6\linewidth}}{
        \cellcolor[rgb]{0.9,0.9,0.9}{
            \makecell[{{p{\linewidth}}}]{
                \texttt{\tiny{[P1$\rangle$GM]}}
                \texttt{MOVE: P, (1,7)} \\
            }
        }
    }
    & & \\ \\

    \theutterance \stepcounter{utterance}  
    & & & \multicolumn{2}{p{0.3\linewidth}}{
        \cellcolor[rgb]{0.9,0.9,0.9}{
            \makecell[{{p{\linewidth}}}]{
                \texttt{\tiny{[GM$|$GM]}}
                \texttt{MOVE: P, (1,7)} \\
            }
        }
    }
    & & \\ \\

    \theutterance \stepcounter{utterance}  
    & & & \multicolumn{2}{p{0.3\linewidth}}{
        \cellcolor[rgb]{0.9,0.9,0.9}{
            \makecell[{{p{\linewidth}}}]{
                \texttt{\tiny{[GM$|$GM]}}
                \texttt{Moved 'P' to (1,7) successfully. Your updated grid looks like this:} \\
\texttt{```} \\
\texttt{ 1234567} \\
\texttt{╔══╤════╗} \\
\texttt{║◌T│D◌◌◌║ 1} \\
\texttt{╟──┘◌◌┌─╢ 2} \\
\texttt{║SA◌◌◌│◌║ 3} \\
\texttt{╟─┐◌┌─┼─╢ 4} \\
\texttt{║◌├─┤◌│O║ 5} \\
\texttt{╟─┤◌├─┤◌║ 6} \\
\texttt{║P│◌│M│◌║ 7} \\
\texttt{╚═╧═╧═╧═╝} \\
\texttt{```} \\
            }
        }
    }
    & & \\ \\

    \theutterance \stepcounter{utterance}  
    & & \multicolumn{4}{p{0.6\linewidth}}{
        \cellcolor[rgb]{0.9,0.9,0.9}{
            \makecell[{{p{\linewidth}}}]{
                \texttt{\tiny{[GM$\rangle$P2]}}
                \texttt{Moved 'M' to (7,7) successfully. Your updated grid looks like this:} \\
\texttt{```} \\
\texttt{ 1234567} \\
\texttt{╔══╤════╗} \\
\texttt{║◌S│◌◌◌◌║ 1} \\
\texttt{╟──┘◌◌┌─╢ 2} \\
\texttt{║A◌O◌◌│◌║ 3} \\
\texttt{╟─┐◌┌─┼─╢ 4} \\
\texttt{║◌├─┤T│D║ 5} \\
\texttt{╟─┤◌├─┤◌║ 6} \\
\texttt{║P│◌│◌│M║ 7} \\
\texttt{╚═╧═╧═╧═╝} \\
\texttt{```} \\
\texttt{You are currently playing round 3 of maximum 28.} \\
\texttt{You have collectively accumulated 2 of 16 penalties.} \\
\texttt{The other player moved an object on their grid.} \\
\\ 
\\ 
\texttt{What is your next command?} \\
            }
        }
    }
    & \\ \\

    \theutterance \stepcounter{utterance}  
    & & & \multicolumn{4}{p{0.6\linewidth}}{
        \cellcolor[rgb]{0.9,0.9,0.9}{
            \makecell[{{p{\linewidth}}}]{
                \texttt{\tiny{[GM$\langle$P2]}}
                \texttt{MOVE: O, (2,7)} \\
            }
        }
    }
     \\ \\

    \theutterance \stepcounter{utterance}  
    & & & \multicolumn{2}{p{0.3\linewidth}}{
        \cellcolor[rgb]{0.9,0.9,0.9}{
            \makecell[{{p{\linewidth}}}]{
                \texttt{\tiny{[GM$|$GM]}}
                \texttt{MOVE: O, (2,7)} \\
            }
        }
    }
    & & \\ \\

    \theutterance \stepcounter{utterance}  
    & & & \multicolumn{2}{p{0.3\linewidth}}{
        \cellcolor[rgb]{0.9,0.9,0.9}{
            \makecell[{{p{\linewidth}}}]{
                \texttt{\tiny{[GM$|$GM]}}
                \texttt{Penalty: (2,7) is not empty, but contains '│'.} \\
\\ 
\texttt{You have collectively accumulated 3 of 16 penalties.} \\
\texttt{Please try again!} \\
            }
        }
    }
    & & \\ \\

    \theutterance \stepcounter{utterance}  
    & & \multicolumn{4}{p{0.6\linewidth}}{
        \cellcolor[rgb]{0.9,0.9,0.9}{
            \makecell[{{p{\linewidth}}}]{
                \texttt{\tiny{[GM$\rangle$P2]}}
                \texttt{Penalty: (2,7) is not empty, but contains '│'.} \\
\\ 
\texttt{You have collectively accumulated 3 of 16 penalties.} \\
\texttt{Please try again!} \\
            }
        }
    }
    & \\ \\

    \theutterance \stepcounter{utterance}  
    & & & \multicolumn{4}{p{0.6\linewidth}}{
        \cellcolor[rgb]{0.9,0.9,0.9}{
            \makecell[{{p{\linewidth}}}]{
                \texttt{\tiny{[GM$\langle$P2]}}
                \texttt{I see that (2,7) is not empty because it contains a vertical line '│'. This means that in row 7, some cells are separated by walls, and (2,7) might not be a valid cell for an object. Looking at the grid, row 7 is drawn as "║P│◌│◌│M║", which indicates that there are vertical lines between columns. Specifically, after column 1 (which has P), there is a vertical line, then column 2 (which is ◌), then another vertical line, then column 3 (which is now ◌ after moving M), then another vertical line, then column 4 (which has M now). But the grid has 7 columns, so this representation might only show columns 1{-}4? However, the top edge has numbers 1 to 7, so there should be 7 columns.} \\
\\ 
\texttt{From the bottom border "╚═╧═╧═╧═╝", it has four "═" segments, which might correspond to four cell groups. This suggests that the grid might have only 4 columns? But the top edge says 1234567, which is confusing.} \\
\\ 
\texttt{To avoid further penalties, I need to reconsider the goal configuration. The other player suggested row 7, but on my board, row 7 seems to have only 4 cells (with vertical lines in between), so it might not be possible to place all 7 objects in row 7. Therefore, I should communicate this issue to the other player and propose a different goal.} \\
\\ 
\texttt{So, my next command should be a SAY command to inform the other player about the grid structure.} \\
\\ 
\texttt{SAY: On my board, row 7 has only 4 cells due to vertical lines. We cannot place all 7 objects there. Let's choose a different row, like row 5 or 6, or use multiple rows. What do you suggest?} \\
            }
        }
    }
     \\ \\

    \theutterance \stepcounter{utterance}  
    & & & \multicolumn{2}{p{0.3\linewidth}}{
        \cellcolor[rgb]{0.9,0.9,0.9}{
            \makecell[{{p{\linewidth}}}]{
                \texttt{\tiny{[GM$|$GM]}}
                \texttt{I see that (2,7) is not empty because it contains a vertical line '│'. This means that in row 7, some cells are separated by walls, and (2,7) might not be a valid cell for an object. Looking at the grid, row 7 is drawn as "║P│◌│◌│M║", which indicates that there are vertical lines between columns. Specifically, after column 1 (which has P), there is a vertical line, then column 2 (which is ◌), then another vertical line, then column 3 (which is now ◌ after moving M), then another vertical line, then column 4 (which has M now). But the grid has 7 columns, so this representation might only show columns 1{-}4? However, the top edge has numbers 1 to 7, so there should be 7 columns.} \\
\\ 
\texttt{From the bottom border "╚═╧═╧═╧═╝", it has four "═" segments, which might correspond to four cell groups. This suggests that the grid might have only 4 columns? But the top edge says 1234567, which is confusing.} \\
\\ 
\texttt{To avoid further penalties, I need to reconsider the goal configuration. The other player suggested row 7, but on my board, row 7 seems to have only 4 cells (with vertical lines in between), so it might not be possible to place all 7 objects in row 7. Therefore, I should communicate this issue to the other player and propose a different goal.} \\
\\ 
\texttt{So, my next command should be a SAY command to inform the other player about the grid structure.} \\
\\ 
\texttt{SAY: On my board, row 7 has only 4 cells due to vertical lines. We cannot place all 7 objects there. Let's choose a different row, like row 5 or 6, or use multiple rows. What do you suggest?} \\
            }
        }
    }
    & & \\ \\

    \theutterance \stepcounter{utterance}  
    & & & \multicolumn{2}{p{0.3\linewidth}}{
        \cellcolor[rgb]{0.9,0.9,0.9}{
            \makecell[{{p{\linewidth}}}]{
                \texttt{\tiny{[GM$|$GM]}}
                \texttt{Invalid format: head is not empty: 'I see that (2,7) is not empty because it contains a vertical line '│'. This means that in row 7, some cells are separated by walls, and (2,7) might not be a valid cell for an object. Looking at the grid, row 7 is drawn as "║P│◌│◌│M║", which indicates that there are vertical lines between columns. Specifically, after column 1 (which has P), there is a vertical line, then column 2 (which is ◌), then another vertical line, then column 3 (which is now ◌ after moving M), then another vertical line, then column 4 (which has M now). But the grid has 7 columns, so this representation might only show columns 1{-}4? However, the top edge has numbers 1 to 7, so there should be 7 columns.} \\
\\ 
\texttt{From the bottom border "╚═╧═╧═╧═╝", it has four "═" segments, which might correspond to four cell groups. This suggests that the grid might have only 4 columns? But the top edge says 1234567, which is confusing.} \\
\\ 
\texttt{To avoid further penalties, I need to reconsider the goal configuration. The other player suggested row 7, but on my board, row 7 seems to have only 4 cells (with vertical lines in between), so it might not be possible to place all 7 objects in row 7. Therefore, I should communicate this issue to the other player and propose a different goal.} \\
\\ 
\texttt{So, my next command should be a SAY command to inform the other player about the grid structure.} \\
\\ 
\texttt{'} \\
            }
        }
    }
    & & \\ \\

    \theutterance \stepcounter{utterance}  
    & & \multicolumn{4}{p{0.6\linewidth}}{
        \cellcolor[rgb]{0.9,0.9,0.9}{
            \makecell[{{p{\linewidth}}}]{
                \texttt{\tiny{[GM$\rangle$P2]}}
                \texttt{Penalty: Your message must not contain anything before the command!} \\
\\ 
\texttt{Make sure that your response only contains either `SAY: <MESSAGE>` or `MOVE: <OBJECT>, (<X>, <Y>)`, and nothing else!} \\
\\ 
\texttt{You have collectively accumulated 4 of 16 penalties.} \\
\texttt{Please try again!} \\
            }
        }
    }
    & \\ \\

    \theutterance \stepcounter{utterance}  
    & & & \multicolumn{4}{p{0.6\linewidth}}{
        \cellcolor[rgb]{0.9,0.9,0.9}{
            \makecell[{{p{\linewidth}}}]{
                \texttt{\tiny{[GM$\langle$P2]}}
                \texttt{SAY: On my board, row 7 has only 4 cells due to vertical lines. We cannot place all 7 objects there. Let's choose a different row, like row 5 or 6, or use multiple rows. What do you suggest?} \\
            }
        }
    }
     \\ \\

    \theutterance \stepcounter{utterance}  
    & & & \multicolumn{2}{p{0.3\linewidth}}{
        \cellcolor[rgb]{0.9,0.9,0.9}{
            \makecell[{{p{\linewidth}}}]{
                \texttt{\tiny{[GM$|$GM]}}
                \texttt{SAY: On my board, row 7 has only 4 cells due to vertical lines. We cannot place all 7 objects there. Let's choose a different row, like row 5 or 6, or use multiple rows. What do you suggest?} \\
            }
        }
    }
    & & \\ \\

    \theutterance \stepcounter{utterance}  
    & & \multicolumn{4}{p{0.6\linewidth}}{
        \cellcolor[rgb]{0.9,0.9,0.9}{
            \makecell[{{p{\linewidth}}}]{
                \texttt{\tiny{[P1$\langle$GM]}}
                \texttt{Moved 'P' to (1,7) successfully. Your updated grid looks like this:} \\
\texttt{```} \\
\texttt{ 1234567} \\
\texttt{╔══╤════╗} \\
\texttt{║◌T│D◌◌◌║ 1} \\
\texttt{╟──┘◌◌┌─╢ 2} \\
\texttt{║SA◌◌◌│◌║ 3} \\
\texttt{╟─┐◌┌─┼─╢ 4} \\
\texttt{║◌├─┤◌│O║ 5} \\
\texttt{╟─┤◌├─┤◌║ 6} \\
\texttt{║P│◌│M│◌║ 7} \\
\texttt{╚═╧═╧═╧═╝} \\
\texttt{```} \\
\texttt{You are currently playing round 3 of maximum 28.} \\
\texttt{You have collectively accumulated 4 of 16 penalties.} \\
\texttt{The other player sent this message:} \\
\\ 
\\ 
\texttt{"On my board, row 7 has only 4 cells due to vertical lines. We cannot place all 7 objects there. Let's choose a different row, like row 5 or 6, or use multiple rows. What do you suggest?"} \\
\\ 
\\ 
\texttt{What is your next command?} \\
            }
        }
    }
    & \\ \\

    \theutterance \stepcounter{utterance}  
    & \multicolumn{4}{p{0.6\linewidth}}{
        \cellcolor[rgb]{0.9,0.9,0.9}{
            \makecell[{{p{\linewidth}}}]{
                \texttt{\tiny{[P1$\rangle$GM]}}
                \texttt{SAY: I see the issue. Row 7 only has 4 cells (columns 1,2,3,4) because of the vertical lines. Let's use row 5 instead {-} it has more space? Alternatively, we can use row 1 which has 7 columns? Checking my grid: row 1 has columns 1{-}7 but has some obstacles. Actually, row 5 has columns 1{-}7? Looking: row 5 has cells at (1,5) to (7,5) but with some lines? Wait, let me check: row 5 has (1,5) empty? (2,5) has a horizontal line? Actually, let's list all empty cells. Perhaps we should use row 3 or row 4? Or agree on specific coordinates for each object. What's your suggestion?} \\
            }
        }
    }
    & & \\ \\

    \theutterance \stepcounter{utterance}  
    & & & \multicolumn{2}{p{0.3\linewidth}}{
        \cellcolor[rgb]{0.9,0.9,0.9}{
            \makecell[{{p{\linewidth}}}]{
                \texttt{\tiny{[GM$|$GM]}}
                \texttt{SAY: I see the issue. Row 7 only has 4 cells (columns 1,2,3,4) because of the vertical lines. Let's use row 5 instead {-} it has more space? Alternatively, we can use row 1 which has 7 columns? Checking my grid: row 1 has columns 1{-}7 but has some obstacles. Actually, row 5 has columns 1{-}7? Looking: row 5 has cells at (1,5) to (7,5) but with some lines? Wait, let me check: row 5 has (1,5) empty? (2,5) has a horizontal line? Actually, let's list all empty cells. Perhaps we should use row 3 or row 4? Or agree on specific coordinates for each object. What's your suggestion?} \\
            }
        }
    }
    & & \\ \\

    \theutterance \stepcounter{utterance}  
    & & \multicolumn{4}{p{0.6\linewidth}}{
        \cellcolor[rgb]{0.9,0.9,0.9}{
            \makecell[{{p{\linewidth}}}]{
                \texttt{\tiny{[GM$\rangle$P2]}}
                \texttt{Your message has been relayed to the other player.} \\
\\ 
\\ 
\texttt{You are currently playing round 4 of maximum 28.} \\
\texttt{You have collectively accumulated 4 of 16 penalties.} \\
\texttt{The other player sent this message:} \\
\\ 
\\ 
\texttt{"I see the issue. Row 7 only has 4 cells (columns 1,2,3,4) because of the vertical lines. Let's use row 5 instead {-} it has more space? Alternatively, we can use row 1 which has 7 columns? Checking my grid: row 1 has columns 1{-}7 but has some obstacles. Actually, row 5 has columns 1{-}7? Looking: row 5 has cells at (1,5) to (7,5) but with some lines? Wait, let me check: row 5 has (1,5) empty? (2,5) has a horizontal line? Actually, let's list all empty cells. Perhaps we should use row 3 or row 4? Or agree on specific coordinates for each object. What's your suggestion?"} \\
\\ 
\\ 
\texttt{What is your next command?} \\
            }
        }
    }
    & \\ \\

    \theutterance \stepcounter{utterance}  
    & & & \multicolumn{4}{p{0.6\linewidth}}{
        \cellcolor[rgb]{0.9,0.9,0.9}{
            \makecell[{{p{\linewidth}}}]{
                \texttt{\tiny{[GM$\langle$P2]}}
                \texttt{SAY: Let's place objects in row 5: P(1,5), O(2,5), T(3,5), S(4,5), D(5,5), A(6,5), M(7,5). Check if row 5 has 7 cells on your board. On mine, it seems row 5 has cells from 1 to 7. I'll verify by moving T to (3,5) first.} \\
            }
        }
    }
     \\ \\

    \theutterance \stepcounter{utterance}  
    & & & \multicolumn{2}{p{0.3\linewidth}}{
        \cellcolor[rgb]{0.9,0.9,0.9}{
            \makecell[{{p{\linewidth}}}]{
                \texttt{\tiny{[GM$|$GM]}}
                \texttt{SAY: Let's place objects in row 5: P(1,5), O(2,5), T(3,5), S(4,5), D(5,5), A(6,5), M(7,5). Check if row 5 has 7 cells on your board. On mine, it seems row 5 has cells from 1 to 7. I'll verify by moving T to (3,5) first.} \\
            }
        }
    }
    & & \\ \\

    \theutterance \stepcounter{utterance}  
    & & \multicolumn{4}{p{0.6\linewidth}}{
        \cellcolor[rgb]{0.9,0.9,0.9}{
            \makecell[{{p{\linewidth}}}]{
                \texttt{\tiny{[P1$\langle$GM]}}
                \texttt{Your message has been relayed to the other player.} \\
\\ 
\\ 
\texttt{You are currently playing round 4 of maximum 28.} \\
\texttt{You have collectively accumulated 4 of 16 penalties.} \\
\texttt{The other player sent this message:} \\
\\ 
\\ 
\texttt{"Let's place objects in row 5: P(1,5), O(2,5), T(3,5), S(4,5), D(5,5), A(6,5), M(7,5). Check if row 5 has 7 cells on your board. On mine, it seems row 5 has cells from 1 to 7. I'll verify by moving T to (3,5) first."} \\
\\ 
\\ 
\texttt{What is your next command?} \\
            }
        }
    }
    & \\ \\

    \theutterance \stepcounter{utterance}  
    & \multicolumn{4}{p{0.6\linewidth}}{
        \cellcolor[rgb]{0.9,0.9,0.9}{
            \makecell[{{p{\linewidth}}}]{
                \texttt{\tiny{[P1$\rangle$GM]}}
                \texttt{SAY: Row 5 on my grid has cells from 1 to 7, but there are some lines and objects already. Currently, (1,5) is empty (◌), (2,5) has a horizontal/vertical junction? Actually, (2,5) is '├─' which is not empty. (3,5) is '┤'? Wait, let's list: (1,5)=◌, (2,5)=├ (part of line), (3,5)=┤ (line), (4,5)=◌, (5,5)=│ (line), (6,5)=O, (7,5)=? Actually, looking at grid: row 5 is "║◌├─┤◌│O║" so cells: (1,5)=◌, (2,5)=├ (not empty), (3,5)=─ (not empty), (4,5)=┤ (not empty), (5,5)=◌, (6,5)=│ (not empty), (7,5)=O. So only (1,5) and (5,5) are empty? We cannot use row 5. Let's use row 1 instead? Row 1: "║◌T│D◌◌◌║" so (1,1)=◌, (2,1)=T, (3,1)=│ (line), (4,1)=D, (5,1)=◌, (6,1)=◌, (7,1)=◌. So only (1,1), (5,1), (6,1), (7,1) are empty. Not enough. Perhaps row 3? Row 3: "║SA◌◌◌│◌║" so (1,3)=S, (2,3)=A, (3,3)=◌, (4,3)=◌, (5,3)=◌, (6,3)=│ (line), (7,3)=◌. So (3,3),(4,3),(5,3),(7,3) empty. Not enough. Row 6? Row 6: "╟─┤◌├─┤◌║" mostly lines. Only (3,6)=◌ and (7,6)=◌? Actually, (1,6)=─ (line), (2,6)=┤ (line), (3,6)=◌, (4,6)=├ (line), (5,6)=─ (line), (6,6)=┤ (line), (7,6)=◌. So only two empty cells. Row 7 has four empty cells: (1,7)=P (I moved P there), (2,7)=◌, (3,7)=M, (4,7)=◌. So we need to use multiple rows. Let's assign: P(1,7), O(2,7), T(3,7), S(4,7), D(1,6), A(3,6), M(7,6)? Or better, use rows 3 and 1? What do you think?} \\
            }
        }
    }
    & & \\ \\

    \theutterance \stepcounter{utterance}  
    & & & \multicolumn{2}{p{0.3\linewidth}}{
        \cellcolor[rgb]{0.9,0.9,0.9}{
            \makecell[{{p{\linewidth}}}]{
                \texttt{\tiny{[GM$|$GM]}}
                \texttt{SAY: Row 5 on my grid has cells from 1 to 7, but there are some lines and objects already. Currently, (1,5) is empty (◌), (2,5) has a horizontal/vertical junction? Actually, (2,5) is '├─' which is not empty. (3,5) is '┤'? Wait, let's list: (1,5)=◌, (2,5)=├ (part of line), (3,5)=┤ (line), (4,5)=◌, (5,5)=│ (line), (6,5)=O, (7,5)=? Actually, looking at grid: row 5 is "║◌├─┤◌│O║" so cells: (1,5)=◌, (2,5)=├ (not empty), (3,5)=─ (not empty), (4,5)=┤ (not empty), (5,5)=◌, (6,5)=│ (not empty), (7,5)=O. So only (1,5) and (5,5) are empty? We cannot use row 5. Let's use row 1 instead? Row 1: "║◌T│D◌◌◌║" so (1,1)=◌, (2,1)=T, (3,1)=│ (line), (4,1)=D, (5,1)=◌, (6,1)=◌, (7,1)=◌. So only (1,1), (5,1), (6,1), (7,1) are empty. Not enough. Perhaps row 3? Row 3: "║SA◌◌◌│◌║" so (1,3)=S, (2,3)=A, (3,3)=◌, (4,3)=◌, (5,3)=◌, (6,3)=│ (line), (7,3)=◌. So (3,3),(4,3),(5,3),(7,3) empty. Not enough. Row 6? Row 6: "╟─┤◌├─┤◌║" mostly lines. Only (3,6)=◌ and (7,6)=◌? Actually, (1,6)=─ (line), (2,6)=┤ (line), (3,6)=◌, (4,6)=├ (line), (5,6)=─ (line), (6,6)=┤ (line), (7,6)=◌. So only two empty cells. Row 7 has four empty cells: (1,7)=P (I moved P there), (2,7)=◌, (3,7)=M, (4,7)=◌. So we need to use multiple rows. Let's assign: P(1,7), O(2,7), T(3,7), S(4,7), D(1,6), A(3,6), M(7,6)? Or better, use rows 3 and 1? What do you think?} \\
            }
        }
    }
    & & \\ \\

    \theutterance \stepcounter{utterance}  
    & & \multicolumn{4}{p{0.6\linewidth}}{
        \cellcolor[rgb]{0.9,0.9,0.9}{
            \makecell[{{p{\linewidth}}}]{
                \texttt{\tiny{[GM$\rangle$P2]}}
                \texttt{Your message has been relayed to the other player.} \\
\\ 
\\ 
\texttt{You are currently playing round 5 of maximum 28.} \\
\texttt{You have collectively accumulated 4 of 16 penalties.} \\
\texttt{The other player sent this message:} \\
\\ 
\\ 
\texttt{"Row 5 on my grid has cells from 1 to 7, but there are some lines and objects already. Currently, (1,5) is empty (◌), (2,5) has a horizontal/vertical junction? Actually, (2,5) is '├─' which is not empty. (3,5) is '┤'? Wait, let's list: (1,5)=◌, (2,5)=├ (part of line), (3,5)=┤ (line), (4,5)=◌, (5,5)=│ (line), (6,5)=O, (7,5)=? Actually, looking at grid: row 5 is "║◌├─┤◌│O║" so cells: (1,5)=◌, (2,5)=├ (not empty), (3,5)=─ (not empty), (4,5)=┤ (not empty), (5,5)=◌, (6,5)=│ (not empty), (7,5)=O. So only (1,5) and (5,5) are empty? We cannot use row 5. Let's use row 1 instead? Row 1: "║◌T│D◌◌◌║" so (1,1)=◌, (2,1)=T, (3,1)=│ (line), (4,1)=D, (5,1)=◌, (6,1)=◌, (7,1)=◌. So only (1,1), (5,1), (6,1), (7,1) are empty. Not enough. Perhaps row 3? Row 3: "║SA◌◌◌│◌║" so (1,3)=S, (2,3)=A, (3,3)=◌, (4,3)=◌, (5,3)=◌, (6,3)=│ (line), (7,3)=◌. So (3,3),(4,3),(5,3),(7,3) empty. Not enough. Row 6? Row 6: "╟─┤◌├─┤◌║" mostly lines. Only (3,6)=◌ and (7,6)=◌? Actually, (1,6)=─ (line), (2,6)=┤ (line), (3,6)=◌, (4,6)=├ (line), (5,6)=─ (line), (6,6)=┤ (line), (7,6)=◌. So only two empty cells. Row 7 has four empty cells: (1,7)=P (I moved P there), (2,7)=◌, (3,7)=M, (4,7)=◌. So we need to use multiple rows. Let's assign: P(1,7), O(2,7), T(3,7), S(4,7), D(1,6), A(3,6), M(7,6)? Or better, use rows 3 and 1? What do you think?"} \\
\\ 
\\ 
\texttt{What is your next command?} \\
            }
        }
    }
    & \\ \\

    \theutterance \stepcounter{utterance}  
    & & & \multicolumn{4}{p{0.6\linewidth}}{
        \cellcolor[rgb]{0.9,0.9,0.9}{
            \makecell[{{p{\linewidth}}}]{
                \texttt{\tiny{[GM$\langle$P2]}}
                \texttt{SAY: Let's use a consistent pattern: place all objects in the leftmost column (column 1) from row 1 to 7: P(1,1), O(1,2), T(1,3), S(1,4), D(1,5), A(1,6), M(1,7). Check if column 1 has 7 cells on your board. On mine, column 1 seems free from row 1 to 7. I'll start by moving A from (1,3) to (1,6).} \\
            }
        }
    }
     \\ \\

    \theutterance \stepcounter{utterance}  
    & & & \multicolumn{2}{p{0.3\linewidth}}{
        \cellcolor[rgb]{0.9,0.9,0.9}{
            \makecell[{{p{\linewidth}}}]{
                \texttt{\tiny{[GM$|$GM]}}
                \texttt{SAY: Let's use a consistent pattern: place all objects in the leftmost column (column 1) from row 1 to 7: P(1,1), O(1,2), T(1,3), S(1,4), D(1,5), A(1,6), M(1,7). Check if column 1 has 7 cells on your board. On mine, column 1 seems free from row 1 to 7. I'll start by moving A from (1,3) to (1,6).} \\
            }
        }
    }
    & & \\ \\

    \theutterance \stepcounter{utterance}  
    & & \multicolumn{4}{p{0.6\linewidth}}{
        \cellcolor[rgb]{0.9,0.9,0.9}{
            \makecell[{{p{\linewidth}}}]{
                \texttt{\tiny{[P1$\langle$GM]}}
                \texttt{Your message has been relayed to the other player.} \\
\\ 
\\ 
\texttt{You are currently playing round 5 of maximum 28.} \\
\texttt{You have collectively accumulated 4 of 16 penalties.} \\
\texttt{The other player sent this message:} \\
\\ 
\\ 
\texttt{"Let's use a consistent pattern: place all objects in the leftmost column (column 1) from row 1 to 7: P(1,1), O(1,2), T(1,3), S(1,4), D(1,5), A(1,6), M(1,7). Check if column 1 has 7 cells on your board. On mine, column 1 seems free from row 1 to 7. I'll start by moving A from (1,3) to (1,6)."} \\
\\ 
\\ 
\texttt{What is your next command?} \\
            }
        }
    }
    & \\ \\

    \theutterance \stepcounter{utterance}  
    & \multicolumn{4}{p{0.6\linewidth}}{
        \cellcolor[rgb]{0.9,0.9,0.9}{
            \makecell[{{p{\linewidth}}}]{
                \texttt{\tiny{[P1$\rangle$GM]}}
                \texttt{SAY: Column 1 on my grid: row1 (1,1)=◌ (empty), row2 (1,2)=╟ (line, not empty), row3 (1,3)=S, row4 (1,4)=╟ (line), row5 (1,5)=◌ (empty), row6 (1,6)=╟ (line), row7 (1,7)=P. So only (1,1) and (1,5) are empty. We cannot use column 1. Let's list all empty cells on my grid to find a common pattern. Empty cells: (1,1), (3,3), (4,3), (5,3), (7,3), (1,5), (5,5), (3,6), (7,6), (2,7), (4,7). So 11 empty cells. We need to assign 7 objects. Let's choose: P(1,7) {[}already{]}, O(7,3), T(1,1), S(3,3), D(4,3), A(5,3), M(7,6). Does that work? Or we can use a diagonal? What empty cells do you have?} \\
            }
        }
    }
    & & \\ \\

    \theutterance \stepcounter{utterance}  
    & & & \multicolumn{2}{p{0.3\linewidth}}{
        \cellcolor[rgb]{0.9,0.9,0.9}{
            \makecell[{{p{\linewidth}}}]{
                \texttt{\tiny{[GM$|$GM]}}
                \texttt{SAY: Column 1 on my grid: row1 (1,1)=◌ (empty), row2 (1,2)=╟ (line, not empty), row3 (1,3)=S, row4 (1,4)=╟ (line), row5 (1,5)=◌ (empty), row6 (1,6)=╟ (line), row7 (1,7)=P. So only (1,1) and (1,5) are empty. We cannot use column 1. Let's list all empty cells on my grid to find a common pattern. Empty cells: (1,1), (3,3), (4,3), (5,3), (7,3), (1,5), (5,5), (3,6), (7,6), (2,7), (4,7). So 11 empty cells. We need to assign 7 objects. Let's choose: P(1,7) {[}already{]}, O(7,3), T(1,1), S(3,3), D(4,3), A(5,3), M(7,6). Does that work? Or we can use a diagonal? What empty cells do you have?} \\
            }
        }
    }
    & & \\ \\

    \theutterance \stepcounter{utterance}  
    & & \multicolumn{4}{p{0.6\linewidth}}{
        \cellcolor[rgb]{0.9,0.9,0.9}{
            \makecell[{{p{\linewidth}}}]{
                \texttt{\tiny{[GM$\rangle$P2]}}
                \texttt{Your message has been relayed to the other player.} \\
\\ 
\\ 
\texttt{You are currently playing round 6 of maximum 28.} \\
\texttt{You have collectively accumulated 4 of 16 penalties.} \\
\texttt{The other player sent this message:} \\
\\ 
\\ 
\texttt{"Column 1 on my grid: row1 (1,1)=◌ (empty), row2 (1,2)=╟ (line, not empty), row3 (1,3)=S, row4 (1,4)=╟ (line), row5 (1,5)=◌ (empty), row6 (1,6)=╟ (line), row7 (1,7)=P. So only (1,1) and (1,5) are empty. We cannot use column 1. Let's list all empty cells on my grid to find a common pattern. Empty cells: (1,1), (3,3), (4,3), (5,3), (7,3), (1,5), (5,5), (3,6), (7,6), (2,7), (4,7). So 11 empty cells. We need to assign 7 objects. Let's choose: P(1,7) {[}already{]}, O(7,3), T(1,1), S(3,3), D(4,3), A(5,3), M(7,6). Does that work? Or we can use a diagonal? What empty cells do you have?"} \\
\\ 
\\ 
\texttt{What is your next command?} \\
            }
        }
    }
    & \\ \\

    \theutterance \stepcounter{utterance}  
    & & & \multicolumn{4}{p{0.6\linewidth}}{
        \cellcolor[rgb]{0.9,0.9,0.9}{
            \makecell[{{p{\linewidth}}}]{
                \texttt{\tiny{[GM$\langle$P2]}}
                \texttt{SAY: My empty cells are: (1,1), (2,1), (4,1), (5,1), (6,1), (7,1), (2,2), (3,2), (4,2), (5,2), (6,2), (7,2), (2,3), (4,3), (5,3), (6,3), (7,3), (1,4), (2,4), (3,4), (4,4), (5,4), (6,4), (7,4), (1,5), (2,5), (3,5), (4,5), (6,5), (7,5), (1,6), (2,6), (4,6), (5,6), (6,6), (7,6), (2,7), (4,7), (5,7), (6,7), (7,7). Let's assign: P(1,7), O(2,7), T(3,7), S(4,7), D(5,7), A(6,7), M(7,7). Row 7 has 7 cells on my board. Confirm if your row 7 has 7 cells? If not, we can use row 1: P(1,1), O(2,1), T(3,1), S(4,1), D(5,1), A(6,1), M(7,1).} \\
            }
        }
    }
     \\ \\

    \theutterance \stepcounter{utterance}  
    & & & \multicolumn{2}{p{0.3\linewidth}}{
        \cellcolor[rgb]{0.9,0.9,0.9}{
            \makecell[{{p{\linewidth}}}]{
                \texttt{\tiny{[GM$|$GM]}}
                \texttt{SAY: My empty cells are: (1,1), (2,1), (4,1), (5,1), (6,1), (7,1), (2,2), (3,2), (4,2), (5,2), (6,2), (7,2), (2,3), (4,3), (5,3), (6,3), (7,3), (1,4), (2,4), (3,4), (4,4), (5,4), (6,4), (7,4), (1,5), (2,5), (3,5), (4,5), (6,5), (7,5), (1,6), (2,6), (4,6), (5,6), (6,6), (7,6), (2,7), (4,7), (5,7), (6,7), (7,7). Let's assign: P(1,7), O(2,7), T(3,7), S(4,7), D(5,7), A(6,7), M(7,7). Row 7 has 7 cells on my board. Confirm if your row 7 has 7 cells? If not, we can use row 1: P(1,1), O(2,1), T(3,1), S(4,1), D(5,1), A(6,1), M(7,1).} \\
            }
        }
    }
    & & \\ \\

    \theutterance \stepcounter{utterance}  
    & & \multicolumn{4}{p{0.6\linewidth}}{
        \cellcolor[rgb]{0.9,0.9,0.9}{
            \makecell[{{p{\linewidth}}}]{
                \texttt{\tiny{[P1$\langle$GM]}}
                \texttt{Your message has been relayed to the other player.} \\
\\ 
\\ 
\texttt{You are currently playing round 6 of maximum 28.} \\
\texttt{You have collectively accumulated 4 of 16 penalties.} \\
\texttt{The other player sent this message:} \\
\\ 
\\ 
\texttt{"My empty cells are: (1,1), (2,1), (4,1), (5,1), (6,1), (7,1), (2,2), (3,2), (4,2), (5,2), (6,2), (7,2), (2,3), (4,3), (5,3), (6,3), (7,3), (1,4), (2,4), (3,4), (4,4), (5,4), (6,4), (7,4), (1,5), (2,5), (3,5), (4,5), (6,5), (7,5), (1,6), (2,6), (4,6), (5,6), (6,6), (7,6), (2,7), (4,7), (5,7), (6,7), (7,7). Let's assign: P(1,7), O(2,7), T(3,7), S(4,7), D(5,7), A(6,7), M(7,7). Row 7 has 7 cells on my board. Confirm if your row 7 has 7 cells? If not, we can use row 1: P(1,1), O(2,1), T(3,1), S(4,1), D(5,1), A(6,1), M(7,1)."} \\
\\ 
\\ 
\texttt{What is your next command?} \\
            }
        }
    }
    & \\ \\

    \theutterance \stepcounter{utterance}  
    & \multicolumn{4}{p{0.6\linewidth}}{
        \cellcolor[rgb]{0.9,0.9,0.9}{
            \makecell[{{p{\linewidth}}}]{
                \texttt{\tiny{[P1$\rangle$GM]}}
                \texttt{SAY: My row 7 only has 4 cells: columns 1,2,3,4. The grid is not a full 7x7; it has walls and lines. My row 1 has 7 columns but only (1,1), (5,1), (6,1), (7,1) are empty. So we cannot use a full row. Let's use the empty cells I have: (1,1), (3,3), (4,3), (5,3), (7,3), (1,5), (5,5), (3,6), (7,6), (2,7), (4,7). I propose: P(1,7) {[}already{]}, O(2,7), T(1,1), S(3,3), D(4,3), A(5,3), M(7,6). This uses 7 distinct empty cells. What are your coordinates for these? If they match your empty cells, we can proceed.} \\
            }
        }
    }
    & & \\ \\

    \theutterance \stepcounter{utterance}  
    & & & \multicolumn{2}{p{0.3\linewidth}}{
        \cellcolor[rgb]{0.9,0.9,0.9}{
            \makecell[{{p{\linewidth}}}]{
                \texttt{\tiny{[GM$|$GM]}}
                \texttt{SAY: My row 7 only has 4 cells: columns 1,2,3,4. The grid is not a full 7x7; it has walls and lines. My row 1 has 7 columns but only (1,1), (5,1), (6,1), (7,1) are empty. So we cannot use a full row. Let's use the empty cells I have: (1,1), (3,3), (4,3), (5,3), (7,3), (1,5), (5,5), (3,6), (7,6), (2,7), (4,7). I propose: P(1,7) {[}already{]}, O(2,7), T(1,1), S(3,3), D(4,3), A(5,3), M(7,6). This uses 7 distinct empty cells. What are your coordinates for these? If they match your empty cells, we can proceed.} \\
            }
        }
    }
    & & \\ \\

    \theutterance \stepcounter{utterance}  
    & & \multicolumn{4}{p{0.6\linewidth}}{
        \cellcolor[rgb]{0.9,0.9,0.9}{
            \makecell[{{p{\linewidth}}}]{
                \texttt{\tiny{[GM$\rangle$P2]}}
                \texttt{Your message has been relayed to the other player.} \\
\\ 
\\ 
\texttt{You are currently playing round 7 of maximum 28.} \\
\texttt{You have collectively accumulated 4 of 16 penalties.} \\
\texttt{The other player sent this message:} \\
\\ 
\\ 
\texttt{"My row 7 only has 4 cells: columns 1,2,3,4. The grid is not a full 7x7; it has walls and lines. My row 1 has 7 columns but only (1,1), (5,1), (6,1), (7,1) are empty. So we cannot use a full row. Let's use the empty cells I have: (1,1), (3,3), (4,3), (5,3), (7,3), (1,5), (5,5), (3,6), (7,6), (2,7), (4,7). I propose: P(1,7) {[}already{]}, O(2,7), T(1,1), S(3,3), D(4,3), A(5,3), M(7,6). This uses 7 distinct empty cells. What are your coordinates for these? If they match your empty cells, we can proceed."} \\
\\ 
\\ 
\texttt{What is your next command?} \\
            }
        }
    }
    & \\ \\

    \theutterance \stepcounter{utterance}  
    & & & \multicolumn{4}{p{0.6\linewidth}}{
        \cellcolor[rgb]{0.9,0.9,0.9}{
            \makecell[{{p{\linewidth}}}]{
                \texttt{\tiny{[GM$\langle$P2]}}
                \texttt{SAY: For your proposal: P(1,7), O(2,7), T(1,1), S(3,3), D(4,3), A(5,3), M(7,6). On my board, (1,7) is P (already), (2,7) is empty, (1,1) is empty, (3,3) is empty, (4,3) is empty, (5,3) is empty, (7,6) is empty. So all are available. I agree! I will start moving objects to these positions.} \\
            }
        }
    }
     \\ \\

    \theutterance \stepcounter{utterance}  
    & & & \multicolumn{2}{p{0.3\linewidth}}{
        \cellcolor[rgb]{0.9,0.9,0.9}{
            \makecell[{{p{\linewidth}}}]{
                \texttt{\tiny{[GM$|$GM]}}
                \texttt{SAY: For your proposal: P(1,7), O(2,7), T(1,1), S(3,3), D(4,3), A(5,3), M(7,6). On my board, (1,7) is P (already), (2,7) is empty, (1,1) is empty, (3,3) is empty, (4,3) is empty, (5,3) is empty, (7,6) is empty. So all are available. I agree! I will start moving objects to these positions.} \\
            }
        }
    }
    & & \\ \\

    \theutterance \stepcounter{utterance}  
    & & \multicolumn{4}{p{0.6\linewidth}}{
        \cellcolor[rgb]{0.9,0.9,0.9}{
            \makecell[{{p{\linewidth}}}]{
                \texttt{\tiny{[P1$\langle$GM]}}
                \texttt{Your message has been relayed to the other player.} \\
\\ 
\\ 
\texttt{You are currently playing round 7 of maximum 28.} \\
\texttt{You have collectively accumulated 4 of 16 penalties.} \\
\texttt{The other player sent this message:} \\
\\ 
\\ 
\texttt{"For your proposal: P(1,7), O(2,7), T(1,1), S(3,3), D(4,3), A(5,3), M(7,6). On my board, (1,7) is P (already), (2,7) is empty, (1,1) is empty, (3,3) is empty, (4,3) is empty, (5,3) is empty, (7,6) is empty. So all are available. I agree! I will start moving objects to these positions."} \\
\\ 
\\ 
\texttt{What is your next command?} \\
            }
        }
    }
    & \\ \\

    \theutterance \stepcounter{utterance}  
    & \multicolumn{4}{p{0.6\linewidth}}{
        \cellcolor[rgb]{0.9,0.9,0.9}{
            \makecell[{{p{\linewidth}}}]{
                \texttt{\tiny{[P1$\rangle$GM]}}
                \texttt{SAY: Great! I will also move objects to these positions. I already have P at (1,7). I'll move O to (2,7) next.} \\
            }
        }
    }
    & & \\ \\

    \theutterance \stepcounter{utterance}  
    & & & \multicolumn{2}{p{0.3\linewidth}}{
        \cellcolor[rgb]{0.9,0.9,0.9}{
            \makecell[{{p{\linewidth}}}]{
                \texttt{\tiny{[GM$|$GM]}}
                \texttt{SAY: Great! I will also move objects to these positions. I already have P at (1,7). I'll move O to (2,7) next.} \\
            }
        }
    }
    & & \\ \\

    \theutterance \stepcounter{utterance}  
    & & \multicolumn{4}{p{0.6\linewidth}}{
        \cellcolor[rgb]{0.9,0.9,0.9}{
            \makecell[{{p{\linewidth}}}]{
                \texttt{\tiny{[GM$\rangle$P2]}}
                \texttt{Your message has been relayed to the other player.} \\
\\ 
\\ 
\texttt{You are currently playing round 8 of maximum 28.} \\
\texttt{You have collectively accumulated 4 of 16 penalties.} \\
\texttt{The other player sent this message:} \\
\\ 
\\ 
\texttt{"Great! I will also move objects to these positions. I already have P at (1,7). I'll move O to (2,7) next."} \\
\\ 
\\ 
\texttt{What is your next command?} \\
            }
        }
    }
    & \\ \\

    \theutterance \stepcounter{utterance}  
    & & & \multicolumn{4}{p{0.6\linewidth}}{
        \cellcolor[rgb]{0.9,0.9,0.9}{
            \makecell[{{p{\linewidth}}}]{
                \texttt{\tiny{[GM$\langle$P2]}}
                \texttt{MOVE: O, (2,7)} \\
            }
        }
    }
     \\ \\

    \theutterance \stepcounter{utterance}  
    & & & \multicolumn{2}{p{0.3\linewidth}}{
        \cellcolor[rgb]{0.9,0.9,0.9}{
            \makecell[{{p{\linewidth}}}]{
                \texttt{\tiny{[GM$|$GM]}}
                \texttt{MOVE: O, (2,7)} \\
            }
        }
    }
    & & \\ \\

    \theutterance \stepcounter{utterance}  
    & & & \multicolumn{2}{p{0.3\linewidth}}{
        \cellcolor[rgb]{0.9,0.9,0.9}{
            \makecell[{{p{\linewidth}}}]{
                \texttt{\tiny{[GM$|$GM]}}
                \texttt{Penalty: (2,7) is not empty, but contains '│'.} \\
\\ 
\texttt{You have collectively accumulated 5 of 16 penalties.} \\
\texttt{Please try again!} \\
            }
        }
    }
    & & \\ \\

    \theutterance \stepcounter{utterance}  
    & & \multicolumn{4}{p{0.6\linewidth}}{
        \cellcolor[rgb]{0.9,0.9,0.9}{
            \makecell[{{p{\linewidth}}}]{
                \texttt{\tiny{[GM$\rangle$P2]}}
                \texttt{Penalty: (2,7) is not empty, but contains '│'.} \\
\\ 
\texttt{You have collectively accumulated 5 of 16 penalties.} \\
\texttt{Please try again!} \\
            }
        }
    }
    & \\ \\

    \theutterance \stepcounter{utterance}  
    & & & \multicolumn{4}{p{0.6\linewidth}}{
        \cellcolor[rgb]{0.9,0.9,0.9}{
            \makecell[{{p{\linewidth}}}]{
                \texttt{\tiny{[GM$\langle$P2]}}
                \texttt{SAY: On my board, (2,7) has a vertical line '│' and is not empty. We need to adjust the plan. Let's use (4,7) for O instead? Or choose another empty cell. My available cells in row 7 are (2,7) is blocked, but (4,7) is empty. So propose: O(4,7). What do you think?} \\
            }
        }
    }
     \\ \\

    \theutterance \stepcounter{utterance}  
    & & & \multicolumn{2}{p{0.3\linewidth}}{
        \cellcolor[rgb]{0.9,0.9,0.9}{
            \makecell[{{p{\linewidth}}}]{
                \texttt{\tiny{[GM$|$GM]}}
                \texttt{SAY: On my board, (2,7) has a vertical line '│' and is not empty. We need to adjust the plan. Let's use (4,7) for O instead? Or choose another empty cell. My available cells in row 7 are (2,7) is blocked, but (4,7) is empty. So propose: O(4,7). What do you think?} \\
            }
        }
    }
    & & \\ \\

    \theutterance \stepcounter{utterance}  
    & & \multicolumn{4}{p{0.6\linewidth}}{
        \cellcolor[rgb]{0.9,0.9,0.9}{
            \makecell[{{p{\linewidth}}}]{
                \texttt{\tiny{[P1$\langle$GM]}}
                \texttt{Your message has been relayed to the other player.} \\
\\ 
\\ 
\texttt{You are currently playing round 8 of maximum 28.} \\
\texttt{You have collectively accumulated 5 of 16 penalties.} \\
\texttt{The other player sent this message:} \\
\\ 
\\ 
\texttt{"On my board, (2,7) has a vertical line '│' and is not empty. We need to adjust the plan. Let's use (4,7) for O instead? Or choose another empty cell. My available cells in row 7 are (2,7) is blocked, but (4,7) is empty. So propose: O(4,7). What do you think?"} \\
\\ 
\\ 
\texttt{What is your next command?} \\
            }
        }
    }
    & \\ \\

    \theutterance \stepcounter{utterance}  
    & \multicolumn{4}{p{0.6\linewidth}}{
        \cellcolor[rgb]{0.9,0.9,0.9}{
            \makecell[{{p{\linewidth}}}]{
                \texttt{\tiny{[P1$\rangle$GM]}}
                \texttt{SAY: On my board, (2,7) is empty (I checked: it's ◌). But if (2,7) is blocked for you, let's use (4,7) for O. So updated: P(1,7), O(4,7), T(1,1), S(3,3), D(4,3), A(5,3), M(7,6). Does that work?} \\
            }
        }
    }
    & & \\ \\

    \theutterance \stepcounter{utterance}  
    & & & \multicolumn{2}{p{0.3\linewidth}}{
        \cellcolor[rgb]{0.9,0.9,0.9}{
            \makecell[{{p{\linewidth}}}]{
                \texttt{\tiny{[GM$|$GM]}}
                \texttt{SAY: On my board, (2,7) is empty (I checked: it's ◌). But if (2,7) is blocked for you, let's use (4,7) for O. So updated: P(1,7), O(4,7), T(1,1), S(3,3), D(4,3), A(5,3), M(7,6). Does that work?} \\
            }
        }
    }
    & & \\ \\

    \theutterance \stepcounter{utterance}  
    & & \multicolumn{4}{p{0.6\linewidth}}{
        \cellcolor[rgb]{0.9,0.9,0.9}{
            \makecell[{{p{\linewidth}}}]{
                \texttt{\tiny{[GM$\rangle$P2]}}
                \texttt{Your message has been relayed to the other player.} \\
\\ 
\\ 
\texttt{You are currently playing round 9 of maximum 28.} \\
\texttt{You have collectively accumulated 5 of 16 penalties.} \\
\texttt{The other player sent this message:} \\
\\ 
\\ 
\texttt{"On my board, (2,7) is empty (I checked: it's ◌). But if (2,7) is blocked for you, let's use (4,7) for O. So updated: P(1,7), O(4,7), T(1,1), S(3,3), D(4,3), A(5,3), M(7,6). Does that work?"} \\
\\ 
\\ 
\texttt{What is your next command?} \\
            }
        }
    }
    & \\ \\

    \theutterance \stepcounter{utterance}  
    & & & \multicolumn{4}{p{0.6\linewidth}}{
        \cellcolor[rgb]{0.9,0.9,0.9}{
            \makecell[{{p{\linewidth}}}]{
                \texttt{\tiny{[GM$\langle$P2]}}
                \texttt{SAY: Yes, that works! O(4,7) is empty on my board. I will move O to (4,7) now.} \\
            }
        }
    }
     \\ \\

    \theutterance \stepcounter{utterance}  
    & & & \multicolumn{2}{p{0.3\linewidth}}{
        \cellcolor[rgb]{0.9,0.9,0.9}{
            \makecell[{{p{\linewidth}}}]{
                \texttt{\tiny{[GM$|$GM]}}
                \texttt{SAY: Yes, that works! O(4,7) is empty on my board. I will move O to (4,7) now.} \\
            }
        }
    }
    & & \\ \\

    \theutterance \stepcounter{utterance}  
    & & \multicolumn{4}{p{0.6\linewidth}}{
        \cellcolor[rgb]{0.9,0.9,0.9}{
            \makecell[{{p{\linewidth}}}]{
                \texttt{\tiny{[P1$\langle$GM]}}
                \texttt{Your message has been relayed to the other player.} \\
\\ 
\\ 
\texttt{You are currently playing round 9 of maximum 28.} \\
\texttt{You have collectively accumulated 5 of 16 penalties.} \\
\texttt{The other player sent this message:} \\
\\ 
\\ 
\texttt{"Yes, that works! O(4,7) is empty on my board. I will move O to (4,7) now."} \\
\\ 
\\ 
\texttt{What is your next command?} \\
            }
        }
    }
    & \\ \\

    \theutterance \stepcounter{utterance}  
    & \multicolumn{4}{p{0.6\linewidth}}{
        \cellcolor[rgb]{0.9,0.9,0.9}{
            \makecell[{{p{\linewidth}}}]{
                \texttt{\tiny{[P1$\rangle$GM]}}
                \texttt{MOVE: O, (4,7)} \\
            }
        }
    }
    & & \\ \\

    \theutterance \stepcounter{utterance}  
    & & & \multicolumn{2}{p{0.3\linewidth}}{
        \cellcolor[rgb]{0.9,0.9,0.9}{
            \makecell[{{p{\linewidth}}}]{
                \texttt{\tiny{[GM$|$GM]}}
                \texttt{MOVE: O, (4,7)} \\
            }
        }
    }
    & & \\ \\

    \theutterance \stepcounter{utterance}  
    & & & \multicolumn{2}{p{0.3\linewidth}}{
        \cellcolor[rgb]{0.9,0.9,0.9}{
            \makecell[{{p{\linewidth}}}]{
                \texttt{\tiny{[GM$|$GM]}}
                \texttt{Penalty: (4,7) is not empty, but contains '│'.} \\
\\ 
\texttt{You have collectively accumulated 6 of 16 penalties.} \\
\texttt{Please try again!} \\
            }
        }
    }
    & & \\ \\

    \theutterance \stepcounter{utterance}  
    & & \multicolumn{4}{p{0.6\linewidth}}{
        \cellcolor[rgb]{0.9,0.9,0.9}{
            \makecell[{{p{\linewidth}}}]{
                \texttt{\tiny{[P1$\langle$GM]}}
                \texttt{Penalty: (4,7) is not empty, but contains '│'.} \\
\\ 
\texttt{You have collectively accumulated 6 of 16 penalties.} \\
\texttt{Please try again!} \\
            }
        }
    }
    & \\ \\

    \theutterance \stepcounter{utterance}  
    & \multicolumn{4}{p{0.6\linewidth}}{
        \cellcolor[rgb]{0.9,0.9,0.9}{
            \makecell[{{p{\linewidth}}}]{
                \texttt{\tiny{[P1$\rangle$GM]}}
                \texttt{SAY: On my board, (4,7) has a vertical line '│' and is not empty. Let's choose another cell. How about O(7,3)? It is empty on my grid. So new proposal: P(1,7), O(7,3), T(1,1), S(3,3), D(4,3), A(5,3), M(7,6). Please confirm if (7,3) is empty for you.} \\
            }
        }
    }
    & & \\ \\

    \theutterance \stepcounter{utterance}  
    & & & \multicolumn{2}{p{0.3\linewidth}}{
        \cellcolor[rgb]{0.9,0.9,0.9}{
            \makecell[{{p{\linewidth}}}]{
                \texttt{\tiny{[GM$|$GM]}}
                \texttt{SAY: On my board, (4,7) has a vertical line '│' and is not empty. Let's choose another cell. How about O(7,3)? It is empty on my grid. So new proposal: P(1,7), O(7,3), T(1,1), S(3,3), D(4,3), A(5,3), M(7,6). Please confirm if (7,3) is empty for you.} \\
            }
        }
    }
    & & \\ \\

    \theutterance \stepcounter{utterance}  
    & & \multicolumn{4}{p{0.6\linewidth}}{
        \cellcolor[rgb]{0.9,0.9,0.9}{
            \makecell[{{p{\linewidth}}}]{
                \texttt{\tiny{[GM$\rangle$P2]}}
                \texttt{Your message has been relayed to the other player.} \\
\\ 
\\ 
\texttt{You are currently playing round 10 of maximum 28.} \\
\texttt{You have collectively accumulated 6 of 16 penalties.} \\
\texttt{The other player sent this message:} \\
\\ 
\\ 
\texttt{"On my board, (4,7) has a vertical line '│' and is not empty. Let's choose another cell. How about O(7,3)? It is empty on my grid. So new proposal: P(1,7), O(7,3), T(1,1), S(3,3), D(4,3), A(5,3), M(7,6). Please confirm if (7,3) is empty for you."} \\
\\ 
\\ 
\texttt{What is your next command?} \\
            }
        }
    }
    & \\ \\

    \theutterance \stepcounter{utterance}  
    & & & \multicolumn{4}{p{0.6\linewidth}}{
        \cellcolor[rgb]{0.9,0.9,0.9}{
            \makecell[{{p{\linewidth}}}]{
                \texttt{\tiny{[GM$\langle$P2]}}
                \texttt{SAY: (7,3) is empty on my board. So I agree: P(1,7), O(7,3), T(1,1), S(3,3), D(4,3), A(5,3), M(7,6). I will move O to (7,3) first.} \\
            }
        }
    }
     \\ \\

    \theutterance \stepcounter{utterance}  
    & & & \multicolumn{2}{p{0.3\linewidth}}{
        \cellcolor[rgb]{0.9,0.9,0.9}{
            \makecell[{{p{\linewidth}}}]{
                \texttt{\tiny{[GM$|$GM]}}
                \texttt{SAY: (7,3) is empty on my board. So I agree: P(1,7), O(7,3), T(1,1), S(3,3), D(4,3), A(5,3), M(7,6). I will move O to (7,3) first.} \\
            }
        }
    }
    & & \\ \\

    \theutterance \stepcounter{utterance}  
    & & \multicolumn{4}{p{0.6\linewidth}}{
        \cellcolor[rgb]{0.9,0.9,0.9}{
            \makecell[{{p{\linewidth}}}]{
                \texttt{\tiny{[P1$\langle$GM]}}
                \texttt{Your message has been relayed to the other player.} \\
\\ 
\\ 
\texttt{You are currently playing round 10 of maximum 28.} \\
\texttt{You have collectively accumulated 6 of 16 penalties.} \\
\texttt{The other player sent this message:} \\
\\ 
\\ 
\texttt{"(7,3) is empty on my board. So I agree: P(1,7), O(7,3), T(1,1), S(3,3), D(4,3), A(5,3), M(7,6). I will move O to (7,3) first."} \\
\\ 
\\ 
\texttt{What is your next command?} \\
            }
        }
    }
    & \\ \\

    \theutterance \stepcounter{utterance}  
    & \multicolumn{4}{p{0.6\linewidth}}{
        \cellcolor[rgb]{0.9,0.9,0.9}{
            \makecell[{{p{\linewidth}}}]{
                \texttt{\tiny{[P1$\rangle$GM]}}
                \texttt{MOVE: O, (7,3)} \\
            }
        }
    }
    & & \\ \\

    \theutterance \stepcounter{utterance}  
    & & & \multicolumn{2}{p{0.3\linewidth}}{
        \cellcolor[rgb]{0.9,0.9,0.9}{
            \makecell[{{p{\linewidth}}}]{
                \texttt{\tiny{[GM$|$GM]}}
                \texttt{MOVE: O, (7,3)} \\
            }
        }
    }
    & & \\ \\

    \theutterance \stepcounter{utterance}  
    & & & \multicolumn{2}{p{0.3\linewidth}}{
        \cellcolor[rgb]{0.9,0.9,0.9}{
            \makecell[{{p{\linewidth}}}]{
                \texttt{\tiny{[GM$|$GM]}}
                \texttt{Moved 'O' to (7,3) successfully. Your updated grid looks like this:} \\
\texttt{```} \\
\texttt{ 1234567} \\
\texttt{╔══╤════╗} \\
\texttt{║◌T│D◌◌◌║ 1} \\
\texttt{╟──┘◌◌┌─╢ 2} \\
\texttt{║SA◌◌◌│O║ 3} \\
\texttt{╟─┐◌┌─┼─╢ 4} \\
\texttt{║◌├─┤◌│◌║ 5} \\
\texttt{╟─┤◌├─┤◌║ 6} \\
\texttt{║P│◌│M│◌║ 7} \\
\texttt{╚═╧═╧═╧═╝} \\
\texttt{```} \\
            }
        }
    }
    & & \\ \\

    \theutterance \stepcounter{utterance}  
    & & \multicolumn{4}{p{0.6\linewidth}}{
        \cellcolor[rgb]{0.9,0.9,0.9}{
            \makecell[{{p{\linewidth}}}]{
                \texttt{\tiny{[GM$\rangle$P2]}}
                \texttt{Your message has been relayed to the other player.} \\
\\ 
\\ 
\texttt{You are currently playing round 11 of maximum 28.} \\
\texttt{You have collectively accumulated 6 of 16 penalties.} \\
\texttt{The other player moved an object on their grid.} \\
\\ 
\\ 
\texttt{What is your next command?} \\
            }
        }
    }
    & \\ \\

    \theutterance \stepcounter{utterance}  
    & & & \multicolumn{4}{p{0.6\linewidth}}{
        \cellcolor[rgb]{0.9,0.9,0.9}{
            \makecell[{{p{\linewidth}}}]{
                \texttt{\tiny{[GM$\langle$P2]}}
                \texttt{MOVE: O, (7,3)} \\
            }
        }
    }
     \\ \\

    \theutterance \stepcounter{utterance}  
    & & & \multicolumn{2}{p{0.3\linewidth}}{
        \cellcolor[rgb]{0.9,0.9,0.9}{
            \makecell[{{p{\linewidth}}}]{
                \texttt{\tiny{[GM$|$GM]}}
                \texttt{MOVE: O, (7,3)} \\
            }
        }
    }
    & & \\ \\

    \theutterance \stepcounter{utterance}  
    & & & \multicolumn{2}{p{0.3\linewidth}}{
        \cellcolor[rgb]{0.9,0.9,0.9}{
            \makecell[{{p{\linewidth}}}]{
                \texttt{\tiny{[GM$|$GM]}}
                \texttt{Moved 'O' to (7,3) successfully. Your updated grid looks like this:} \\
\texttt{```} \\
\texttt{ 1234567} \\
\texttt{╔══╤════╗} \\
\texttt{║◌S│◌◌◌◌║ 1} \\
\texttt{╟──┘◌◌┌─╢ 2} \\
\texttt{║A◌◌◌◌│O║ 3} \\
\texttt{╟─┐◌┌─┼─╢ 4} \\
\texttt{║◌├─┤T│D║ 5} \\
\texttt{╟─┤◌├─┤◌║ 6} \\
\texttt{║P│◌│◌│M║ 7} \\
\texttt{╚═╧═╧═╧═╝} \\
\texttt{```} \\
            }
        }
    }
    & & \\ \\

    \theutterance \stepcounter{utterance}  
    & & \multicolumn{4}{p{0.6\linewidth}}{
        \cellcolor[rgb]{0.9,0.9,0.9}{
            \makecell[{{p{\linewidth}}}]{
                \texttt{\tiny{[P1$\langle$GM]}}
                \texttt{Moved 'O' to (7,3) successfully. Your updated grid looks like this:} \\
\texttt{```} \\
\texttt{ 1234567} \\
\texttt{╔══╤════╗} \\
\texttt{║◌T│D◌◌◌║ 1} \\
\texttt{╟──┘◌◌┌─╢ 2} \\
\texttt{║SA◌◌◌│O║ 3} \\
\texttt{╟─┐◌┌─┼─╢ 4} \\
\texttt{║◌├─┤◌│◌║ 5} \\
\texttt{╟─┤◌├─┤◌║ 6} \\
\texttt{║P│◌│M│◌║ 7} \\
\texttt{╚═╧═╧═╧═╝} \\
\texttt{```} \\
\texttt{You are currently playing round 11 of maximum 28.} \\
\texttt{You have collectively accumulated 6 of 16 penalties.} \\
\texttt{The other player moved an object on their grid.} \\
\\ 
\\ 
\texttt{What is your next command?} \\
            }
        }
    }
    & \\ \\

    \theutterance \stepcounter{utterance}  
    & \multicolumn{4}{p{0.6\linewidth}}{
        \cellcolor[rgb]{0.9,0.9,0.9}{
            \makecell[{{p{\linewidth}}}]{
                \texttt{\tiny{[P1$\rangle$GM]}}
                \texttt{MOVE: T, (1,1)} \\
            }
        }
    }
    & & \\ \\

    \theutterance \stepcounter{utterance}  
    & & & \multicolumn{2}{p{0.3\linewidth}}{
        \cellcolor[rgb]{0.9,0.9,0.9}{
            \makecell[{{p{\linewidth}}}]{
                \texttt{\tiny{[GM$|$GM]}}
                \texttt{MOVE: T, (1,1)} \\
            }
        }
    }
    & & \\ \\

    \theutterance \stepcounter{utterance}  
    & & & \multicolumn{2}{p{0.3\linewidth}}{
        \cellcolor[rgb]{0.9,0.9,0.9}{
            \makecell[{{p{\linewidth}}}]{
                \texttt{\tiny{[GM$|$GM]}}
                \texttt{Moved 'T' to (1,1) successfully. Your updated grid looks like this:} \\
\texttt{```} \\
\texttt{ 1234567} \\
\texttt{╔══╤════╗} \\
\texttt{║T◌│D◌◌◌║ 1} \\
\texttt{╟──┘◌◌┌─╢ 2} \\
\texttt{║SA◌◌◌│O║ 3} \\
\texttt{╟─┐◌┌─┼─╢ 4} \\
\texttt{║◌├─┤◌│◌║ 5} \\
\texttt{╟─┤◌├─┤◌║ 6} \\
\texttt{║P│◌│M│◌║ 7} \\
\texttt{╚═╧═╧═╧═╝} \\
\texttt{```} \\
            }
        }
    }
    & & \\ \\

    \theutterance \stepcounter{utterance}  
    & & \multicolumn{4}{p{0.6\linewidth}}{
        \cellcolor[rgb]{0.9,0.9,0.9}{
            \makecell[{{p{\linewidth}}}]{
                \texttt{\tiny{[GM$\rangle$P2]}}
                \texttt{Moved 'O' to (7,3) successfully. Your updated grid looks like this:} \\
\texttt{```} \\
\texttt{ 1234567} \\
\texttt{╔══╤════╗} \\
\texttt{║◌S│◌◌◌◌║ 1} \\
\texttt{╟──┘◌◌┌─╢ 2} \\
\texttt{║A◌◌◌◌│O║ 3} \\
\texttt{╟─┐◌┌─┼─╢ 4} \\
\texttt{║◌├─┤T│D║ 5} \\
\texttt{╟─┤◌├─┤◌║ 6} \\
\texttt{║P│◌│◌│M║ 7} \\
\texttt{╚═╧═╧═╧═╝} \\
\texttt{```} \\
\texttt{You are currently playing round 12 of maximum 28.} \\
\texttt{You have collectively accumulated 6 of 16 penalties.} \\
\texttt{The other player moved an object on their grid.} \\
\\ 
\\ 
\texttt{What is your next command?} \\
            }
        }
    }
    & \\ \\

    \theutterance \stepcounter{utterance}  
    & & & \multicolumn{4}{p{0.6\linewidth}}{
        \cellcolor[rgb]{0.9,0.9,0.9}{
            \makecell[{{p{\linewidth}}}]{
                \texttt{\tiny{[GM$\langle$P2]}}
                \texttt{MOVE: T, (1,1)} \\
            }
        }
    }
     \\ \\

    \theutterance \stepcounter{utterance}  
    & & & \multicolumn{2}{p{0.3\linewidth}}{
        \cellcolor[rgb]{0.9,0.9,0.9}{
            \makecell[{{p{\linewidth}}}]{
                \texttt{\tiny{[GM$|$GM]}}
                \texttt{MOVE: T, (1,1)} \\
            }
        }
    }
    & & \\ \\

    \theutterance \stepcounter{utterance}  
    & & & \multicolumn{2}{p{0.3\linewidth}}{
        \cellcolor[rgb]{0.9,0.9,0.9}{
            \makecell[{{p{\linewidth}}}]{
                \texttt{\tiny{[GM$|$GM]}}
                \texttt{Moved 'T' to (1,1) successfully. Your updated grid looks like this:} \\
\texttt{```} \\
\texttt{ 1234567} \\
\texttt{╔══╤════╗} \\
\texttt{║TS│◌◌◌◌║ 1} \\
\texttt{╟──┘◌◌┌─╢ 2} \\
\texttt{║A◌◌◌◌│O║ 3} \\
\texttt{╟─┐◌┌─┼─╢ 4} \\
\texttt{║◌├─┤◌│D║ 5} \\
\texttt{╟─┤◌├─┤◌║ 6} \\
\texttt{║P│◌│◌│M║ 7} \\
\texttt{╚═╧═╧═╧═╝} \\
\texttt{```} \\
            }
        }
    }
    & & \\ \\

    \theutterance \stepcounter{utterance}  
    & & \multicolumn{4}{p{0.6\linewidth}}{
        \cellcolor[rgb]{0.9,0.9,0.9}{
            \makecell[{{p{\linewidth}}}]{
                \texttt{\tiny{[P1$\langle$GM]}}
                \texttt{Moved 'T' to (1,1) successfully. Your updated grid looks like this:} \\
\texttt{```} \\
\texttt{ 1234567} \\
\texttt{╔══╤════╗} \\
\texttt{║T◌│D◌◌◌║ 1} \\
\texttt{╟──┘◌◌┌─╢ 2} \\
\texttt{║SA◌◌◌│O║ 3} \\
\texttt{╟─┐◌┌─┼─╢ 4} \\
\texttt{║◌├─┤◌│◌║ 5} \\
\texttt{╟─┤◌├─┤◌║ 6} \\
\texttt{║P│◌│M│◌║ 7} \\
\texttt{╚═╧═╧═╧═╝} \\
\texttt{```} \\
\texttt{You are currently playing round 12 of maximum 28.} \\
\texttt{You have collectively accumulated 6 of 16 penalties.} \\
\texttt{The other player moved an object on their grid.} \\
\\ 
\\ 
\texttt{What is your next command?} \\
            }
        }
    }
    & \\ \\

    \theutterance \stepcounter{utterance}  
    & \multicolumn{4}{p{0.6\linewidth}}{
        \cellcolor[rgb]{0.9,0.9,0.9}{
            \makecell[{{p{\linewidth}}}]{
                \texttt{\tiny{[P1$\rangle$GM]}}
                \texttt{MOVE: S, (3,3)} \\
            }
        }
    }
    & & \\ \\

    \theutterance \stepcounter{utterance}  
    & & & \multicolumn{2}{p{0.3\linewidth}}{
        \cellcolor[rgb]{0.9,0.9,0.9}{
            \makecell[{{p{\linewidth}}}]{
                \texttt{\tiny{[GM$|$GM]}}
                \texttt{MOVE: S, (3,3)} \\
            }
        }
    }
    & & \\ \\

    \theutterance \stepcounter{utterance}  
    & & & \multicolumn{2}{p{0.3\linewidth}}{
        \cellcolor[rgb]{0.9,0.9,0.9}{
            \makecell[{{p{\linewidth}}}]{
                \texttt{\tiny{[GM$|$GM]}}
                \texttt{Moved 'S' to (3,3) successfully. Your updated grid looks like this:} \\
\texttt{```} \\
\texttt{ 1234567} \\
\texttt{╔══╤════╗} \\
\texttt{║T◌│D◌◌◌║ 1} \\
\texttt{╟──┘◌◌┌─╢ 2} \\
\texttt{║◌AS◌◌│O║ 3} \\
\texttt{╟─┐◌┌─┼─╢ 4} \\
\texttt{║◌├─┤◌│◌║ 5} \\
\texttt{╟─┤◌├─┤◌║ 6} \\
\texttt{║P│◌│M│◌║ 7} \\
\texttt{╚═╧═╧═╧═╝} \\
\texttt{```} \\
            }
        }
    }
    & & \\ \\

    \theutterance \stepcounter{utterance}  
    & & \multicolumn{4}{p{0.6\linewidth}}{
        \cellcolor[rgb]{0.9,0.9,0.9}{
            \makecell[{{p{\linewidth}}}]{
                \texttt{\tiny{[GM$\rangle$P2]}}
                \texttt{Moved 'T' to (1,1) successfully. Your updated grid looks like this:} \\
\texttt{```} \\
\texttt{ 1234567} \\
\texttt{╔══╤════╗} \\
\texttt{║TS│◌◌◌◌║ 1} \\
\texttt{╟──┘◌◌┌─╢ 2} \\
\texttt{║A◌◌◌◌│O║ 3} \\
\texttt{╟─┐◌┌─┼─╢ 4} \\
\texttt{║◌├─┤◌│D║ 5} \\
\texttt{╟─┤◌├─┤◌║ 6} \\
\texttt{║P│◌│◌│M║ 7} \\
\texttt{╚═╧═╧═╧═╝} \\
\texttt{```} \\
\texttt{You are currently playing round 13 of maximum 28.} \\
\texttt{You have collectively accumulated 6 of 16 penalties.} \\
\texttt{The other player moved an object on their grid.} \\
\\ 
\\ 
\texttt{What is your next command?} \\
            }
        }
    }
    & \\ \\

    \theutterance \stepcounter{utterance}  
    & & & \multicolumn{4}{p{0.6\linewidth}}{
        \cellcolor[rgb]{0.9,0.9,0.9}{
            \makecell[{{p{\linewidth}}}]{
                \texttt{\tiny{[GM$\langle$P2]}}
                \texttt{MOVE: S, (3,3)} \\
            }
        }
    }
     \\ \\

    \theutterance \stepcounter{utterance}  
    & & & \multicolumn{2}{p{0.3\linewidth}}{
        \cellcolor[rgb]{0.9,0.9,0.9}{
            \makecell[{{p{\linewidth}}}]{
                \texttt{\tiny{[GM$|$GM]}}
                \texttt{MOVE: S, (3,3)} \\
            }
        }
    }
    & & \\ \\

    \theutterance \stepcounter{utterance}  
    & & & \multicolumn{2}{p{0.3\linewidth}}{
        \cellcolor[rgb]{0.9,0.9,0.9}{
            \makecell[{{p{\linewidth}}}]{
                \texttt{\tiny{[GM$|$GM]}}
                \texttt{Moved 'S' to (3,3) successfully. Your updated grid looks like this:} \\
\texttt{```} \\
\texttt{ 1234567} \\
\texttt{╔══╤════╗} \\
\texttt{║T◌│◌◌◌◌║ 1} \\
\texttt{╟──┘◌◌┌─╢ 2} \\
\texttt{║A◌S◌◌│O║ 3} \\
\texttt{╟─┐◌┌─┼─╢ 4} \\
\texttt{║◌├─┤◌│D║ 5} \\
\texttt{╟─┤◌├─┤◌║ 6} \\
\texttt{║P│◌│◌│M║ 7} \\
\texttt{╚═╧═╧═╧═╝} \\
\texttt{```} \\
            }
        }
    }
    & & \\ \\

    \theutterance \stepcounter{utterance}  
    & & \multicolumn{4}{p{0.6\linewidth}}{
        \cellcolor[rgb]{0.9,0.9,0.9}{
            \makecell[{{p{\linewidth}}}]{
                \texttt{\tiny{[P1$\langle$GM]}}
                \texttt{Moved 'S' to (3,3) successfully. Your updated grid looks like this:} \\
\texttt{```} \\
\texttt{ 1234567} \\
\texttt{╔══╤════╗} \\
\texttt{║T◌│D◌◌◌║ 1} \\
\texttt{╟──┘◌◌┌─╢ 2} \\
\texttt{║◌AS◌◌│O║ 3} \\
\texttt{╟─┐◌┌─┼─╢ 4} \\
\texttt{║◌├─┤◌│◌║ 5} \\
\texttt{╟─┤◌├─┤◌║ 6} \\
\texttt{║P│◌│M│◌║ 7} \\
\texttt{╚═╧═╧═╧═╝} \\
\texttt{```} \\
\texttt{You are currently playing round 13 of maximum 28.} \\
\texttt{You have collectively accumulated 6 of 16 penalties.} \\
\texttt{The other player moved an object on their grid.} \\
\\ 
\\ 
\texttt{What is your next command?} \\
            }
        }
    }
    & \\ \\

    \theutterance \stepcounter{utterance}  
    & \multicolumn{4}{p{0.6\linewidth}}{
        \cellcolor[rgb]{0.9,0.9,0.9}{
            \makecell[{{p{\linewidth}}}]{
                \texttt{\tiny{[P1$\rangle$GM]}}
                \texttt{MOVE: D, (4,3)} \\
            }
        }
    }
    & & \\ \\

    \theutterance \stepcounter{utterance}  
    & & & \multicolumn{2}{p{0.3\linewidth}}{
        \cellcolor[rgb]{0.9,0.9,0.9}{
            \makecell[{{p{\linewidth}}}]{
                \texttt{\tiny{[GM$|$GM]}}
                \texttt{MOVE: D, (4,3)} \\
            }
        }
    }
    & & \\ \\

    \theutterance \stepcounter{utterance}  
    & & & \multicolumn{2}{p{0.3\linewidth}}{
        \cellcolor[rgb]{0.9,0.9,0.9}{
            \makecell[{{p{\linewidth}}}]{
                \texttt{\tiny{[GM$|$GM]}}
                \texttt{Moved 'D' to (4,3) successfully. Your updated grid looks like this:} \\
\texttt{```} \\
\texttt{ 1234567} \\
\texttt{╔══╤════╗} \\
\texttt{║T◌│◌◌◌◌║ 1} \\
\texttt{╟──┘◌◌┌─╢ 2} \\
\texttt{║◌ASD◌│O║ 3} \\
\texttt{╟─┐◌┌─┼─╢ 4} \\
\texttt{║◌├─┤◌│◌║ 5} \\
\texttt{╟─┤◌├─┤◌║ 6} \\
\texttt{║P│◌│M│◌║ 7} \\
\texttt{╚═╧═╧═╧═╝} \\
\texttt{```} \\
            }
        }
    }
    & & \\ \\

    \theutterance \stepcounter{utterance}  
    & & \multicolumn{4}{p{0.6\linewidth}}{
        \cellcolor[rgb]{0.9,0.9,0.9}{
            \makecell[{{p{\linewidth}}}]{
                \texttt{\tiny{[GM$\rangle$P2]}}
                \texttt{Moved 'S' to (3,3) successfully. Your updated grid looks like this:} \\
\texttt{```} \\
\texttt{ 1234567} \\
\texttt{╔══╤════╗} \\
\texttt{║T◌│◌◌◌◌║ 1} \\
\texttt{╟──┘◌◌┌─╢ 2} \\
\texttt{║A◌S◌◌│O║ 3} \\
\texttt{╟─┐◌┌─┼─╢ 4} \\
\texttt{║◌├─┤◌│D║ 5} \\
\texttt{╟─┤◌├─┤◌║ 6} \\
\texttt{║P│◌│◌│M║ 7} \\
\texttt{╚═╧═╧═╧═╝} \\
\texttt{```} \\
\texttt{You are currently playing round 14 of maximum 28.} \\
\texttt{You have collectively accumulated 6 of 16 penalties.} \\
\texttt{The other player moved an object on their grid.} \\
\\ 
\\ 
\texttt{What is your next command?} \\
            }
        }
    }
    & \\ \\

    \theutterance \stepcounter{utterance}  
    & & & \multicolumn{4}{p{0.6\linewidth}}{
        \cellcolor[rgb]{0.9,0.9,0.9}{
            \makecell[{{p{\linewidth}}}]{
                \texttt{\tiny{[GM$\langle$P2]}}
                \texttt{MOVE: D, (4,3)} \\
            }
        }
    }
     \\ \\

    \theutterance \stepcounter{utterance}  
    & & & \multicolumn{2}{p{0.3\linewidth}}{
        \cellcolor[rgb]{0.9,0.9,0.9}{
            \makecell[{{p{\linewidth}}}]{
                \texttt{\tiny{[GM$|$GM]}}
                \texttt{MOVE: D, (4,3)} \\
            }
        }
    }
    & & \\ \\

    \theutterance \stepcounter{utterance}  
    & & & \multicolumn{2}{p{0.3\linewidth}}{
        \cellcolor[rgb]{0.9,0.9,0.9}{
            \makecell[{{p{\linewidth}}}]{
                \texttt{\tiny{[GM$|$GM]}}
                \texttt{Moved 'D' to (4,3) successfully. Your updated grid looks like this:} \\
\texttt{```} \\
\texttt{ 1234567} \\
\texttt{╔══╤════╗} \\
\texttt{║T◌│◌◌◌◌║ 1} \\
\texttt{╟──┘◌◌┌─╢ 2} \\
\texttt{║A◌SD◌│O║ 3} \\
\texttt{╟─┐◌┌─┼─╢ 4} \\
\texttt{║◌├─┤◌│◌║ 5} \\
\texttt{╟─┤◌├─┤◌║ 6} \\
\texttt{║P│◌│◌│M║ 7} \\
\texttt{╚═╧═╧═╧═╝} \\
\texttt{```} \\
            }
        }
    }
    & & \\ \\

    \theutterance \stepcounter{utterance}  
    & & \multicolumn{4}{p{0.6\linewidth}}{
        \cellcolor[rgb]{0.9,0.9,0.9}{
            \makecell[{{p{\linewidth}}}]{
                \texttt{\tiny{[P1$\langle$GM]}}
                \texttt{Moved 'D' to (4,3) successfully. Your updated grid looks like this:} \\
\texttt{```} \\
\texttt{ 1234567} \\
\texttt{╔══╤════╗} \\
\texttt{║T◌│◌◌◌◌║ 1} \\
\texttt{╟──┘◌◌┌─╢ 2} \\
\texttt{║◌ASD◌│O║ 3} \\
\texttt{╟─┐◌┌─┼─╢ 4} \\
\texttt{║◌├─┤◌│◌║ 5} \\
\texttt{╟─┤◌├─┤◌║ 6} \\
\texttt{║P│◌│M│◌║ 7} \\
\texttt{╚═╧═╧═╧═╝} \\
\texttt{```} \\
\texttt{You are currently playing round 14 of maximum 28.} \\
\texttt{You have collectively accumulated 6 of 16 penalties.} \\
\texttt{The other player moved an object on their grid.} \\
\\ 
\\ 
\texttt{What is your next command?} \\
            }
        }
    }
    & \\ \\

    \theutterance \stepcounter{utterance}  
    & \multicolumn{4}{p{0.6\linewidth}}{
        \cellcolor[rgb]{0.9,0.9,0.9}{
            \makecell[{{p{\linewidth}}}]{
                \texttt{\tiny{[P1$\rangle$GM]}}
                \texttt{MOVE: A, (5,3)} \\
            }
        }
    }
    & & \\ \\

    \theutterance \stepcounter{utterance}  
    & & & \multicolumn{2}{p{0.3\linewidth}}{
        \cellcolor[rgb]{0.9,0.9,0.9}{
            \makecell[{{p{\linewidth}}}]{
                \texttt{\tiny{[GM$|$GM]}}
                \texttt{MOVE: A, (5,3)} \\
            }
        }
    }
    & & \\ \\

    \theutterance \stepcounter{utterance}  
    & & & \multicolumn{2}{p{0.3\linewidth}}{
        \cellcolor[rgb]{0.9,0.9,0.9}{
            \makecell[{{p{\linewidth}}}]{
                \texttt{\tiny{[GM$|$GM]}}
                \texttt{Moved 'A' to (5,3) successfully. Your updated grid looks like this:} \\
\texttt{```} \\
\texttt{ 1234567} \\
\texttt{╔══╤════╗} \\
\texttt{║T◌│◌◌◌◌║ 1} \\
\texttt{╟──┘◌◌┌─╢ 2} \\
\texttt{║◌◌SDA│O║ 3} \\
\texttt{╟─┐◌┌─┼─╢ 4} \\
\texttt{║◌├─┤◌│◌║ 5} \\
\texttt{╟─┤◌├─┤◌║ 6} \\
\texttt{║P│◌│M│◌║ 7} \\
\texttt{╚═╧═╧═╧═╝} \\
\texttt{```} \\
            }
        }
    }
    & & \\ \\

    \theutterance \stepcounter{utterance}  
    & & \multicolumn{4}{p{0.6\linewidth}}{
        \cellcolor[rgb]{0.9,0.9,0.9}{
            \makecell[{{p{\linewidth}}}]{
                \texttt{\tiny{[GM$\rangle$P2]}}
                \texttt{Moved 'D' to (4,3) successfully. Your updated grid looks like this:} \\
\texttt{```} \\
\texttt{ 1234567} \\
\texttt{╔══╤════╗} \\
\texttt{║T◌│◌◌◌◌║ 1} \\
\texttt{╟──┘◌◌┌─╢ 2} \\
\texttt{║A◌SD◌│O║ 3} \\
\texttt{╟─┐◌┌─┼─╢ 4} \\
\texttt{║◌├─┤◌│◌║ 5} \\
\texttt{╟─┤◌├─┤◌║ 6} \\
\texttt{║P│◌│◌│M║ 7} \\
\texttt{╚═╧═╧═╧═╝} \\
\texttt{```} \\
\texttt{You are currently playing round 15 of maximum 28.} \\
\texttt{You have collectively accumulated 6 of 16 penalties.} \\
\texttt{The other player moved an object on their grid.} \\
\\ 
\\ 
\texttt{What is your next command?} \\
            }
        }
    }
    & \\ \\

    \theutterance \stepcounter{utterance}  
    & & & \multicolumn{4}{p{0.6\linewidth}}{
        \cellcolor[rgb]{0.9,0.9,0.9}{
            \makecell[{{p{\linewidth}}}]{
                \texttt{\tiny{[GM$\langle$P2]}}
                \texttt{MOVE: A, (5,3)} \\
            }
        }
    }
     \\ \\

    \theutterance \stepcounter{utterance}  
    & & & \multicolumn{2}{p{0.3\linewidth}}{
        \cellcolor[rgb]{0.9,0.9,0.9}{
            \makecell[{{p{\linewidth}}}]{
                \texttt{\tiny{[GM$|$GM]}}
                \texttt{MOVE: A, (5,3)} \\
            }
        }
    }
    & & \\ \\

    \theutterance \stepcounter{utterance}  
    & & & \multicolumn{2}{p{0.3\linewidth}}{
        \cellcolor[rgb]{0.9,0.9,0.9}{
            \makecell[{{p{\linewidth}}}]{
                \texttt{\tiny{[GM$|$GM]}}
                \texttt{Moved 'A' to (5,3) successfully. Your updated grid looks like this:} \\
\texttt{```} \\
\texttt{ 1234567} \\
\texttt{╔══╤════╗} \\
\texttt{║T◌│◌◌◌◌║ 1} \\
\texttt{╟──┘◌◌┌─╢ 2} \\
\texttt{║◌◌SDA│O║ 3} \\
\texttt{╟─┐◌┌─┼─╢ 4} \\
\texttt{║◌├─┤◌│◌║ 5} \\
\texttt{╟─┤◌├─┤◌║ 6} \\
\texttt{║P│◌│◌│M║ 7} \\
\texttt{╚═╧═╧═╧═╝} \\
\texttt{```} \\
            }
        }
    }
    & & \\ \\

    \theutterance \stepcounter{utterance}  
    & & \multicolumn{4}{p{0.6\linewidth}}{
        \cellcolor[rgb]{0.9,0.9,0.9}{
            \makecell[{{p{\linewidth}}}]{
                \texttt{\tiny{[P1$\langle$GM]}}
                \texttt{Moved 'A' to (5,3) successfully. Your updated grid looks like this:} \\
\texttt{```} \\
\texttt{ 1234567} \\
\texttt{╔══╤════╗} \\
\texttt{║T◌│◌◌◌◌║ 1} \\
\texttt{╟──┘◌◌┌─╢ 2} \\
\texttt{║◌◌SDA│O║ 3} \\
\texttt{╟─┐◌┌─┼─╢ 4} \\
\texttt{║◌├─┤◌│◌║ 5} \\
\texttt{╟─┤◌├─┤◌║ 6} \\
\texttt{║P│◌│M│◌║ 7} \\
\texttt{╚═╧═╧═╧═╝} \\
\texttt{```} \\
\texttt{You are currently playing round 15 of maximum 28.} \\
\texttt{You have collectively accumulated 6 of 16 penalties.} \\
\texttt{The other player moved an object on their grid.} \\
\\ 
\\ 
\texttt{What is your next command?} \\
            }
        }
    }
    & \\ \\

    \theutterance \stepcounter{utterance}  
    & \multicolumn{4}{p{0.6\linewidth}}{
        \cellcolor[rgb]{0.9,0.9,0.9}{
            \makecell[{{p{\linewidth}}}]{
                \texttt{\tiny{[P1$\rangle$GM]}}
                \texttt{MOVE: M, (7,6)} \\
            }
        }
    }
    & & \\ \\

    \theutterance \stepcounter{utterance}  
    & & & \multicolumn{2}{p{0.3\linewidth}}{
        \cellcolor[rgb]{0.9,0.9,0.9}{
            \makecell[{{p{\linewidth}}}]{
                \texttt{\tiny{[GM$|$GM]}}
                \texttt{MOVE: M, (7,6)} \\
            }
        }
    }
    & & \\ \\

    \theutterance \stepcounter{utterance}  
    & & & \multicolumn{2}{p{0.3\linewidth}}{
        \cellcolor[rgb]{0.9,0.9,0.9}{
            \makecell[{{p{\linewidth}}}]{
                \texttt{\tiny{[GM$|$GM]}}
                \texttt{Moved 'M' to (7,6) successfully. Your updated grid looks like this:} \\
\texttt{```} \\
\texttt{ 1234567} \\
\texttt{╔══╤════╗} \\
\texttt{║T◌│◌◌◌◌║ 1} \\
\texttt{╟──┘◌◌┌─╢ 2} \\
\texttt{║◌◌SDA│O║ 3} \\
\texttt{╟─┐◌┌─┼─╢ 4} \\
\texttt{║◌├─┤◌│◌║ 5} \\
\texttt{╟─┤◌├─┤M║ 6} \\
\texttt{║P│◌│◌│◌║ 7} \\
\texttt{╚═╧═╧═╧═╝} \\
\texttt{```} \\
            }
        }
    }
    & & \\ \\

    \theutterance \stepcounter{utterance}  
    & & \multicolumn{4}{p{0.6\linewidth}}{
        \cellcolor[rgb]{0.9,0.9,0.9}{
            \makecell[{{p{\linewidth}}}]{
                \texttt{\tiny{[GM$\rangle$P2]}}
                \texttt{Moved 'A' to (5,3) successfully. Your updated grid looks like this:} \\
\texttt{```} \\
\texttt{ 1234567} \\
\texttt{╔══╤════╗} \\
\texttt{║T◌│◌◌◌◌║ 1} \\
\texttt{╟──┘◌◌┌─╢ 2} \\
\texttt{║◌◌SDA│O║ 3} \\
\texttt{╟─┐◌┌─┼─╢ 4} \\
\texttt{║◌├─┤◌│◌║ 5} \\
\texttt{╟─┤◌├─┤◌║ 6} \\
\texttt{║P│◌│◌│M║ 7} \\
\texttt{╚═╧═╧═╧═╝} \\
\texttt{```} \\
\texttt{You are currently playing round 16 of maximum 28.} \\
\texttt{You have collectively accumulated 6 of 16 penalties.} \\
\texttt{The other player moved an object on their grid.} \\
\\ 
\\ 
\texttt{What is your next command?} \\
            }
        }
    }
    & \\ \\

    \theutterance \stepcounter{utterance}  
    & & & \multicolumn{4}{p{0.6\linewidth}}{
        \cellcolor[rgb]{0.9,0.9,0.9}{
            \makecell[{{p{\linewidth}}}]{
                \texttt{\tiny{[GM$\langle$P2]}}
                \texttt{MOVE: M, (7,6)} \\
            }
        }
    }
     \\ \\

    \theutterance \stepcounter{utterance}  
    & & & \multicolumn{2}{p{0.3\linewidth}}{
        \cellcolor[rgb]{0.9,0.9,0.9}{
            \makecell[{{p{\linewidth}}}]{
                \texttt{\tiny{[GM$|$GM]}}
                \texttt{MOVE: M, (7,6)} \\
            }
        }
    }
    & & \\ \\

    \theutterance \stepcounter{utterance}  
    & & & \multicolumn{2}{p{0.3\linewidth}}{
        \cellcolor[rgb]{0.9,0.9,0.9}{
            \makecell[{{p{\linewidth}}}]{
                \texttt{\tiny{[GM$|$GM]}}
                \texttt{Moved 'M' to (7,6) successfully. Your updated grid looks like this:} \\
\texttt{```} \\
\texttt{ 1234567} \\
\texttt{╔══╤════╗} \\
\texttt{║T◌│◌◌◌◌║ 1} \\
\texttt{╟──┘◌◌┌─╢ 2} \\
\texttt{║◌◌SDA│O║ 3} \\
\texttt{╟─┐◌┌─┼─╢ 4} \\
\texttt{║◌├─┤◌│◌║ 5} \\
\texttt{╟─┤◌├─┤M║ 6} \\
\texttt{║P│◌│◌│◌║ 7} \\
\texttt{╚═╧═╧═╧═╝} \\
\texttt{```} \\
            }
        }
    }
    & & \\ \\

    \theutterance \stepcounter{utterance}  
    & & \multicolumn{4}{p{0.6\linewidth}}{
        \cellcolor[rgb]{0.9,0.9,0.9}{
            \makecell[{{p{\linewidth}}}]{
                \texttt{\tiny{[P1$\langle$GM]}}
                \texttt{Moved 'M' to (7,6) successfully. Your updated grid looks like this:} \\
\texttt{```} \\
\texttt{ 1234567} \\
\texttt{╔══╤════╗} \\
\texttt{║T◌│◌◌◌◌║ 1} \\
\texttt{╟──┘◌◌┌─╢ 2} \\
\texttt{║◌◌SDA│O║ 3} \\
\texttt{╟─┐◌┌─┼─╢ 4} \\
\texttt{║◌├─┤◌│◌║ 5} \\
\texttt{╟─┤◌├─┤M║ 6} \\
\texttt{║P│◌│◌│◌║ 7} \\
\texttt{╚═╧═╧═╧═╝} \\
\texttt{```} \\
\texttt{You are currently playing round 16 of maximum 28.} \\
\texttt{You have collectively accumulated 6 of 16 penalties.} \\
\texttt{The other player moved an object on their grid.} \\
\\ 
\\ 
\texttt{What is your next command?} \\
            }
        }
    }
    & \\ \\

    \theutterance \stepcounter{utterance}  
    & \multicolumn{4}{p{0.6\linewidth}}{
        \cellcolor[rgb]{0.9,0.9,0.9}{
            \makecell[{{p{\linewidth}}}]{
                \texttt{\tiny{[P1$\rangle$GM]}}
                \texttt{SAY: I have moved all objects to the agreed positions: P(1,7), O(7,3), T(1,1), S(3,3), D(4,3), A(5,3), M(7,6). Are you finished? If so, please say "finished!"} \\
            }
        }
    }
    & & \\ \\

    \theutterance \stepcounter{utterance}  
    & & & \multicolumn{2}{p{0.3\linewidth}}{
        \cellcolor[rgb]{0.9,0.9,0.9}{
            \makecell[{{p{\linewidth}}}]{
                \texttt{\tiny{[GM$|$GM]}}
                \texttt{SAY: I have moved all objects to the agreed positions: P(1,7), O(7,3), T(1,1), S(3,3), D(4,3), A(5,3), M(7,6). Are you finished? If so, please say "finished!"} \\
            }
        }
    }
    & & \\ \\

    \theutterance \stepcounter{utterance}  
    & & \multicolumn{4}{p{0.6\linewidth}}{
        \cellcolor[rgb]{0.9,0.9,0.9}{
            \makecell[{{p{\linewidth}}}]{
                \texttt{\tiny{[GM$\rangle$P2]}}
                \texttt{Moved 'M' to (7,6) successfully. Your updated grid looks like this:} \\
\texttt{```} \\
\texttt{ 1234567} \\
\texttt{╔══╤════╗} \\
\texttt{║T◌│◌◌◌◌║ 1} \\
\texttt{╟──┘◌◌┌─╢ 2} \\
\texttt{║◌◌SDA│O║ 3} \\
\texttt{╟─┐◌┌─┼─╢ 4} \\
\texttt{║◌├─┤◌│◌║ 5} \\
\texttt{╟─┤◌├─┤M║ 6} \\
\texttt{║P│◌│◌│◌║ 7} \\
\texttt{╚═╧═╧═╧═╝} \\
\texttt{```} \\
\texttt{You are currently playing round 17 of maximum 28.} \\
\texttt{You have collectively accumulated 6 of 16 penalties.} \\
\texttt{The other player sent this message:} \\
\\ 
\\ 
\texttt{"I have moved all objects to the agreed positions: P(1,7), O(7,3), T(1,1), S(3,3), D(4,3), A(5,3), M(7,6). Are you finished? If so, please say "finished!""} \\
\\ 
\\ 
\texttt{What is your next command?} \\
            }
        }
    }
    & \\ \\

    \theutterance \stepcounter{utterance}  
    & & & \multicolumn{4}{p{0.6\linewidth}}{
        \cellcolor[rgb]{0.9,0.9,0.9}{
            \makecell[{{p{\linewidth}}}]{
                \texttt{\tiny{[GM$\langle$P2]}}
                \texttt{SAY: finished!} \\
            }
        }
    }
     \\ \\

    \theutterance \stepcounter{utterance}  
    & & & \multicolumn{2}{p{0.3\linewidth}}{
        \cellcolor[rgb]{0.9,0.9,0.9}{
            \makecell[{{p{\linewidth}}}]{
                \texttt{\tiny{[GM$|$GM]}}
                \texttt{SAY: finished!} \\
            }
        }
    }
    & & \\ \\

    \theutterance \stepcounter{utterance}  
    & & & \multicolumn{2}{p{0.3\linewidth}}{
        \cellcolor[rgb]{0.9,0.9,0.9}{
            \makecell[{{p{\linewidth}}}]{
                \texttt{\tiny{[GM$|$GM]}}
                \texttt{true} \\
            }
        }
    }
    & & \\ \\

    \theutterance \stepcounter{utterance}  
    & & & \multicolumn{2}{p{0.3\linewidth}}{
        \cellcolor[rgb]{0.9,0.9,0.9}{
            \makecell[{{p{\linewidth}}}]{
                \texttt{\tiny{[GM$|$GM]}}
                \texttt{* success: True} \\
\texttt{* lose: False} \\
\texttt{* aborted: False} \\
\texttt{{-}{-}{-}{-}{-}{-}{-}} \\
\texttt{* Shifts: 7.00} \\
\texttt{* Max Shifts: 12.00} \\
\texttt{* Min Shifts: 6.00} \\
\texttt{* End Distance Sum: 0.00} \\
\texttt{* Init Distance Sum: 19.71} \\
\texttt{* Expected Distance Sum: 29.33} \\
\texttt{* Penalties: 6.00} \\
\texttt{* Max Penalties: 16.00} \\
\texttt{* Rounds: 17.00} \\
\texttt{* Max Rounds: 28.00} \\
\texttt{* Object Count: 7.00} \\
            }
        }
    }
    & & \\ \\

    \theutterance \stepcounter{utterance}  
    & & & \multicolumn{2}{p{0.3\linewidth}}{
        \cellcolor[rgb]{0.9,0.9,0.9}{
            \makecell[{{p{\linewidth}}}]{
                \texttt{\tiny{[GM$|$GM]}}
                \texttt{Initial Grids:} \\
\texttt{```} \\
\texttt{Player 1:     Player 2:} \\
\texttt{ 1234567       1234567} \\
\texttt{╔══╤════╗     ╔══╤════╗} \\
\texttt{║◌T│D◌◌◌║ 1   ║◌S│◌◌◌◌║ 1} \\
\texttt{╟──┘◌◌┌─╢ 2   ╟──┘◌◌┌─╢ 2} \\
\texttt{║SA◌◌◌│◌║ 3   ║A◌O◌◌│◌║ 3} \\
\texttt{╟─┐◌┌─┼─╢ 4   ╟─┐◌┌─┼─╢ 4} \\
\texttt{║◌├─┤◌│O║ 5   ║◌├─┤T│D║ 5} \\
\texttt{╟─┤◌├─┤◌║ 6   ╟─┤◌├─┤◌║ 6} \\
\texttt{║◌│P│M│◌║ 7   ║P│◌│M│◌║ 7} \\
\texttt{╚═╧═╧═╧═╝     ╚═╧═╧═╧═╝} \\
\texttt{```} \\
\\ 
\texttt{Final Grids:} \\
\texttt{```} \\
\texttt{Player 1:     Player 2:} \\
\texttt{ 1234567       1234567} \\
\texttt{╔══╤════╗     ╔══╤════╗} \\
\texttt{║T◌│◌◌◌◌║ 1   ║T◌│◌◌◌◌║ 1} \\
\texttt{╟──┘◌◌┌─╢ 2   ╟──┘◌◌┌─╢ 2} \\
\texttt{║◌◌SDA│O║ 3   ║◌◌SDA│O║ 3} \\
\texttt{╟─┐◌┌─┼─╢ 4   ╟─┐◌┌─┼─╢ 4} \\
\texttt{║◌├─┤◌│◌║ 5   ║◌├─┤◌│◌║ 5} \\
\texttt{╟─┤◌├─┤M║ 6   ╟─┤◌├─┤M║ 6} \\
\texttt{║P│◌│◌│◌║ 7   ║P│◌│◌│◌║ 7} \\
\texttt{╚═╧═╧═╧═╝     ╚═╧═╧═╧═╝} \\
\texttt{```} \\
            }
        }
    }
    & & \\ \\

\end{supertabular}
}

\end{document}
